% concatenated from main.tex
\documentclass[12pt]{amsart}
\usepackage{amsmath,amscd,amsbsy,amssymb,amsthm,amsfonts}
\usepackage{latexsym,url,bm}
\usepackage{cite,enumitem}
\usepackage{fullpage}
\usepackage{color}
\usepackage[colorlinks=true]{hyperref}
%\usepackage{showkeys}
\usepackage{chngcntr}
\usepackage{apptools}
\usepackage{graphicx}

\newtheorem{claim}{Claim}[section]
\newtheorem{theorem}{Theorem}[section]
\newtheorem{proposition}[theorem]{Proposition}
\newtheorem{lemma}[theorem]{Lemma}
\newtheorem{corollary}[theorem]{Corollary}
%\theoremstyle{remark}
\newtheorem{remark}{Remark}[section]
\newtheorem{example}[remark]{Example}
\theoremstyle{definition}
\newtheorem{definition}[theorem]{Definition}

\setcounter{tocdepth}{1}
\allowdisplaybreaks

\numberwithin{equation}{section}
\renewcommand{\baselinestretch}{1.2}
\newcommand{\p}{\partial}
\newcommand{\e}{\epsilon}
\newcommand{\wt}{\widetilde}
\newcommand{\R}{\mathbb{R}}
\newcommand{\N}{\mathbb{N}}
\newcommand{\Z}{\mathbb{Z}}
\newcommand{\C}{\mathbb{C}}
\newcommand{\F}{\mathcal{F}}
\renewcommand{\H}{{\mathcal H }}
\newcommand{\al}{\alpha}
\newcommand{\f}{\frac}
\newcommand{\n}{\nabla}
\newcommand{\tha}{\theta}
\newcommand{\la}{\lambda}
\newcommand{\bw}{\mathbf{W}}
\newcommand{\bu}{\mathbf{u}}
\newcommand{\nP}{{\mathbf P}}
\newcommand{\uu}{{\underline u}}

\newcommand{\intt}{{\int_0^T}}
\newcommand{\intx}{{\int_{-\pi}^\pi}}
\newcommand{\inty}{{\int_{-h}^{\eta(t,x)}}}

\newcommand{\red}[1]{{\color{red}#1}}
\newcommand{\blue}[1]{{\color{blue}#1}}

\newcommand{\doth}[1]{\Dot{\mathcal{H}}^{#1}}


\keywords{Muskat problem, Elasticity, Well-posedness}
\subjclass[2020]{35R35, 35Q35, 35A01, 35A02}
\pagestyle{plain}

\author{Lizhe Wan}
\address{Beijing International Center for Mathematical Research, Peking University}
\curraddr{}
\email{wanlizhe@pku.edu.cn}

\author{Jiaqi Yang}
\address{School of Mathematics and Statistics, Northwestern Polytechnical University}
\curraddr{}
\email{yjqmath@nwpu.edu.cn, yjqmath@163.com}


\begin{document}

\title{On the well-posedness of two-dimensional Muskat problem with an elastic interface}

\begin{abstract}
We investigate the two-dimensional Muskat problem with a nonlinear elastic interface, for both one-phase and two-phase scenarios.
Following the framework developed by Nguyen \cite{MR4131404, MR4348695}, we demonstrate that the problem is locally well-posed in $H^s$ for $s\geq 2$ for arbitrary initial data.
Furthermore, for the one-phase case and the stable two-phase case $(\rho^+ \leq \rho^-)$, we establish global well-posedness for small initial data in $H^s$ when $s> \f32$.
\end{abstract}

\maketitle
\tableofcontents

\section{Introduction} \label{s:Intro}
Poroelasticity theory, which rigorously couples solid deformation with fluid flow in porous materials, has evolved from a specialized soil consolidation model into a foundational theoretical framework for solving critical challenges across engineering, earth sciences, and biomechanics, see Cheng \cite{cheng2016poroelasticity}. 
Its significance lies in its ability to quantify the bidirectional interaction whereby mechanical loads induce fluid pressure changes and drive flow, while fluid injection or extraction, in turn, deforms the solid matrix. 

This paper studies the coupling between an elastic sheet and a porous medium, a model referred to as the \textbf{Muskat problem with an elastic interface.}
For related models such as hydroelastic waves, the Peskin problem, and the Muskat problem for viscoelastic filtration, one can refer to  Cameron and Strain \cite{MR4673875}, Gahn \cite{MR4959949}, Meirmanov \cite{MR2863467}, and Plotnikov and Toland \cite{MR2812947} for these results.
Our model integrates Darcy’s law for fluids with the Cosserat shell theory (under Kirchhoff’s hypotheses) to describe the elastic sheet. 

A closely related class of models concerns hydroelastic waves, which describes the interaction between elastic structures and hydrodynamic forcing; see \cite{MR2812947}. 
In this setting, the fluid is governed by the incompressible Euler equations.
For developments in numerical simulation, experimental studies, and applications of hydroelastic waves, we refer to Părău \textit{et al.} \cite{MR2812939}. 
The local well-posedness of two-dimensional hydroelastic waves was established by Ambrose and Siegel \cite{MR3656704} and by Liu and Ambrose \cite{MR3608168}. 
Results on local well-posedness for hydroelastic waves with vorticity in arbitrary spatial dimensions can be found in the work of the second author and Wang \cite{MR4104949}. 
Recently, authors obtained a low-regularity well-posedness result for two-dimensional hydroelastic waves in \cite{WanY}.

Let us denote the interface between the fluids (or the fluid and the air) at time $t$ by $\Sigma_t$.
In this paper, we assume throughout that $\Sigma_t$ can be represented by the graph of a time-dependent function $\eta(t,x)$, so that
\begin{equation*}
    \Sigma_t = \{ (x, \eta(t,x)): x\in \R \}.
\end{equation*}
Along the interface $\Sigma_t$, a thin layer of elastic sheet separates the fluid domain into the upper and lower region, denoted by $\Omega^+_t$, $\Omega_t^-$ respectively.
They are given by
\begin{align*}
& \Omega_t^+ = \{ (x, y)\in \R^2: \eta(t,x)<y< \underline{b}^+(x) \}, \\
& \Omega_t^- = \{ (x, y)\in \R^2: \underline{b}^-(x)<y< \eta(t,x) \},
\end{align*}
where $\underline{b}^+(x)$ and $\underline{b}^-(x)$ are parameterizations of the upper and the lower part of the fluid boundary respectively:
\begin{equation*}
    \Gamma^\pm = \{ (x, \underline{b}^\pm(x)): x\in \R \}.
\end{equation*}
\begin{figure}[htbp] 
    \centering 
\includegraphics[width=0.7\textwidth]{Muskat.png}
\caption{Two-phase Muskat problem with an elastic interface} 
    \label{fig}
\end{figure}

The velocities $u^\pm$ and pressure $p^\pm$ of incompressible fluid in each region is governed by Darcy's law:
\begin{equation} \label{Darcy}
\mu^\pm u^\pm + \nabla_{x,y}p^{\pm} = -\rho^{\pm}g \textbf{e}_y, \quad \nabla_{x,y}\cdot u^\pm = 0 \quad \text{in } \Omega^\pm_t, 
\end{equation}
where $g \geq 0$ denotes the gravitational acceleration, $\rho^\pm$ are the densities of the fluids in $\Omega^\pm$, and $\mu^\pm$ are the corresponding viscosity coefficients.

Along the interface $\Sigma_t$, we assume that the normal component of the velocity is continuous, namely,
\begin{equation} \label{NormalVelocity}
    u^+ \cdot n = u^- \cdot n \quad \text{on } \Sigma_t, 
\end{equation}
where $n = \frac{1}{\sqrt{1+\eta_x^2}}(-\eta_x, 1)$ denotes the upward-pointing unit normal vector to $\Sigma_t$.
Then the kinematic boundary condition on $\Sigma_t$ takes the form
\begin{equation} \label{kinematic}
    \p_t \eta = \sqrt{1+\eta_x^2} u^- \cdot n\big|_{\Sigma_t}.
\end{equation}
The dynamic boundary condition on $\Sigma_t$ asserts that the pressure jump across $\Sigma_t$ is balanced by the restoring force generated by the elastic sheet.
\begin{equation} \label{dynamic}
   p^- - p^+ = \sigma \mathbf{E}(\eta) \quad \text{on } \Sigma_t, 
\end{equation}
where $\sigma>0$ is the coefficient of flexural rigidity.
The nonlinear elastic operator $\mathbf{E}(\eta)$, arising from the bending energy of the sheet, is given by
\begin{equation} \label{Elastic}
\begin{aligned}
    \mathbf{E}(\eta) :=&\frac{1}{\sqrt{1+\eta_x^2}} 
    \left[ \frac{1}{\sqrt{1+\eta_x^2}} 
    \left( \frac{\eta_{xx}}{(1+\eta_x^2)^{3/2}} \right)_x \right]_x
    + \frac{1}{2} \left( \frac{\eta_{xx}}{(1+\eta_x^2)^{3/2}} \right)^3\\
    =&\left( \dfrac{1}{1+\eta_x^2}\left(\dfrac{\eta_x}{\sqrt{1+\eta_x^2}}\right)_x\right)_{xx}+\f52\left(\f{\eta_x\eta_{xx}^2}{(1+\eta^2_x)^{\f72}}\right)_x.
\end{aligned}
\end{equation}

On the rigid upper and lower boundary, we assume the non-penetration condition,
\begin{equation} \label{NoPenetration}
    u^\pm \cdot (-\underline{b}^\pm_x, 1) = 0, \quad \text{on } \Gamma^\pm.
\end{equation}

The system \eqref{Darcy}-\eqref{NoPenetration} is referred to as the  \textbf{two-phase} Muskat problem with an elastic interface.
When the upper phase corresponds to air or vacuum, so that $\mu^+ = \rho^+ = p^+=0$, the problem reduces to the \textbf{one-phase} Muskat problem.

\subsection{Formulation of the Muskat problem}
The Muskat problem admits several equivalent formulations; see \cite{MR4520423,MR3415681,MR3171344,MR2313156,MR2993754,MR3608884,MR3071395}. 
In the absence of surface tension, C\'ordoba \& Gancedo \cite{MR2318314} introduced a contour dynamics approach for the infinite-depth Muskat problem without viscosity jump, and proved local well-posedness in $H^3$ for graph interfaces. 
This approach was later extended by C\'ordoba, C\'ordoba \& Gancedo \cite{MR2753607,MR3071395} to include a viscosity jump and to non-graph interfaces satisfying the arc-chord and Rayleigh–Taylor conditions.
Using an arbitrary Lagrangian–Eulerian method, Cheng, Granero \& Shkoller \cite{MR3415681} established local well-posedness for the one-phase problem with a flat bottom, assuming initial surfaces $\eta \in H^2$.
Matioc \cite{MR3861893} later refined the regularity requirement to $\eta \in H^s$ with $s > \frac32$ for the constant-viscosity, infinite-depth case. 
Alazard \& Lazar \cite{MR4097324} provided an alternative proof via paralinearization of the contour dynamics formulation.
Nguyen \& Pausader \cite{MR4090462} reformulated the problem using the Dirichlet–Neumann operator and proved local well-posedness for large data in the critical Sobolev spaces. 
For results on the Muskat problem with surface tension, we refer to \cite{MR4131404,MR2162781,MR3841857}, and references therein.
Concerning global well-posedness, relevant works include \cite{MR3869383,MR3415681,MR4363243,MR3899970,MR3595492,MR2998834,Lazar,MR4348695,MR4655356}.

In this paper, we follow the formulation in Nguyen \cite{MR4131404} and Nguyen \& Pausader \cite{MR4090462}, and rewrite the Muskat problem in terms of the Dirichlet-Neumann operators $G^\pm(\eta)$ associated to $\Omega^\pm_t$. 
For a fixed time $t$, and a given function $f$ defined on $\Sigma$, we let $\phi^\pm$ be the solution of 
\begin{equation*}
\begin{cases}
 \Delta_{x,y} \phi^\pm = 0 \quad \text{in }\Omega^\pm, \\
\phi^\pm = f \quad \text{on }\Sigma, \\
\frac{\p \phi^\pm}{\p\nu^\pm} = 0 \quad \text{on }\Gamma^\pm.
\end{cases}
\end{equation*}
The Dirichlet-Neumann operators $G^\pm(\eta)$ is given by
\begin{equation*}
 G^\pm(\eta) f : = \sqrt{1+\eta_x^2} \frac{\p \phi^\pm}{\p n}.  
\end{equation*}

Then the two-dimensional Muskat problem with an elastic interface can be reformulated using the following result.

\begin{proposition}[\hspace{1sp}\cite{MR4131404}] \label{t:Formuation}
$(i)$ If $(u,p, \eta)$ is a solution to the one-phase Muskat problem, then $\eta$ solves the differential equation
\begin{equation} \label{OneMuskat}
    \p_t \eta = -\frac{1}{\mu^-}G^-(\eta)(\sigma \mathbf{E}(\eta)+ \rho^- g \eta).
\end{equation}
On the other hand, if $\eta$ is a solution of the differential equation \eqref{OneMuskat}, then the one-phase Muskat problem has a solution, in which $\eta$ parameterizes the free surface $\Sigma_t$.

$(ii)$ If $(u^\pm, p^\pm, \eta)$ solve the two-phase Muskat problem if and only if
\begin{equation} \label{TwoMuskatOne}
    \p_t \eta = -\frac{1}{\mu^-}G^-(\eta)f^-,
\end{equation}
where $f^\pm := p^\pm|_{\Sigma_t} + \rho^\pm \eta$ satisfy 
\begin{equation}  \label{TwoMuskatTwo}
\begin{cases}
f^- - f^+ = \sigma \mathbf{E}(\eta) + g(\rho^- - \rho^+)\eta, \\
\frac{1}{\mu^+}G^+(\eta)f^+ = \frac{1}{\mu^-}G^-(\eta)f^-.
\end{cases}
\end{equation}
On the other hand, if $\eta$ is a solution of \eqref{TwoMuskatOne} where $f^\pm$ solve \eqref{TwoMuskatTwo}, then the two-phase Muskat problem has a solution, in which $\eta$ parameterizes the free surface $\Sigma_t$.
\end{proposition}

Proposition \ref{t:Formuation} was proved in Appendix $B$ in \cite{MR4131404} for the Muskat problem with surface tension.
Replacing the capillary terms with elastic terms yields the corresponding result stated above.

The Muskat problem with an elastic interface is essentially a fifth-order quasilinear parabolic equation.
Indeed, at the leading order, one may intuitively approximate $\sigma \mathbf{E}(\eta) + \rho^- g\eta\approx \sigma  \Delta^2 \eta$, so that the one-phase Muskat problem \eqref{OneMuskat} takes the schematic form
\begin{equation} \label{ToyModel}
     \p_t \eta = -\frac{\sigma}{\mu^-}G^-(\eta)\Delta^2 \eta + \text{remainder terms}.
\end{equation}
In the bottomless case where $\Gamma^\pm = \emptyset$, it $\eta(t,x)$ solves \eqref{ToyModel}, so is 
\begin{equation*}
    \eta_\lambda (t,x) = \lambda^{-1}\eta(\lambda^5 t, \lambda x), \quad \forall \lambda >0.
\end{equation*}
Hence, the scale-invariant Sobolev space is $\dot{H}^\f32(\R)$.
This critical regularity coincides with that of the gravity Muskat problem and the Muskat problem with surface tension.
As in the case of water waves with surface tension \cite{MR2805065} and the Muskat problem with surface tension \cite{MR4131404}, the presence of rigid boundaries $\Gamma^\pm$ affects only the low-frequency behavior and does not alter the local well-posedness theory.

\subsection{Main results}
Before stating the main results of this paper, we first introduce the functional spaces used throughout the paper.
We define
\begin{align*}
    &\dot{W}^{1,\infty} : = \{u\in L^1_{loc}(\R) : \nabla u \in L^\infty(\R)  \}, \\
    & Z^s(T) : = C([0,T]; H^s(\R)) \cap L^2([0,T]; H^{s+\f52}(\R)).
\end{align*}

Our first main result is on the local well-posedness of two-dimensional Muskat problem for large data in $H^s$ for $s \geq 2$.
\begin{theorem} \label{t:MainOne}
Let $s\geq 2$, and let $\sigma>0$, $g\geq 0$ be fixed parameters.

$(i)$ \textbf{Local well-posedness for the one-phase problem}: Let $\mu^- >0$, and $\rho^->0$.
We consider either $\Gamma^- = \emptyset$ or $\underline{b}^- \in \dot{W}^{1,\infty}$.
Let $\eta_0 \in H^s(\R) ($\text{or} $H^s(\mathbb{T}))$   satisfy 
\begin{equation*}
    \text{dist}(\eta_0, \Gamma^-)>2h>0.
\end{equation*}
Then there exists a time $T>0$, depending only on $\|\eta_0 \|_{H^s}$, $h$, $s$, $\frac{\rho^- g}{\mu^-}$ and $\frac{\sigma}{\mu^-}$, and a unique solution $\eta \in Z^s(T)$ to the equation \eqref{OneMuskat} such that
\begin{equation*}
    \eta(0,x) = \eta_0, \quad \inf_{t\in[0,T]}\text{dist}(\eta(t), \Gamma^-)>h.
\end{equation*}
Moreover, if $\eta_1$ and $\eta_2$ are two solutions of \eqref{OneMuskat}, then 
\begin{equation} \label{StabilityOne}
    \|\eta_1 - \eta_2 \|_{Z^s(T)} \lesssim_{\|\eta_1\|_{Z^s(T)} + \|\eta_1\|_{Z^s(T)}} \|(\eta_1 - \eta_2)(0,\cdot) \|_{H^s}.
\end{equation}

$(ii)$ \textbf{Local well-posedness for the two-phase problem}: Let $\mu^\pm >0$, and $\rho^\pm>0$.
We consider either $\Gamma^\pm = \emptyset$ or $\underline{b}^\pm \in \dot{W}^{1,\infty}$.
Let $\eta_0 \in H^s(\R) ($\text{or} $H^s(\mathbb{T}))$ satisfy 
\begin{equation*}
    \text{dist}(\eta_0, \Gamma^\pm)>2h>0.
\end{equation*}
Then there exists a time $T>0$, depending only on $\|\eta_0 \|_{H^s}$, $h$, $s$, $\mu^\pm$, $\sigma$ and $g(\rho^- -\rho^+)$, and a unique solution $\eta \in Z^s(T)$ to  \eqref{TwoMuskatOne}-\eqref{TwoMuskatTwo} such that
\begin{equation*}
    \eta(0,x) = \eta_0, \quad \inf_{t\in[0,T]}\text{dist}(\eta(t), \Gamma^\pm)>h.
\end{equation*}
Moreover, if $\eta_1$ and $\eta_2$ are two solutions of \eqref{TwoMuskatOne}-\eqref{TwoMuskatTwo}, then 
\begin{equation} \label{StabilityTwo}
    \|\eta_1 - \eta_2 \|_{Z^s(T)} \lesssim_{\|\eta_1\|_{Z^s(T)} + \|\eta_2\|_{Z^s(T)}} \|(\eta_1 - \eta_2)(0,\cdot) \|_{H^s}.
\end{equation}
\end{theorem}

We make the following remarks about the local well-posedness result.
\begin{remark}
i) Theorem \ref{t:MainOne} establishes well-posedness in the sense of Hadamard, namely: existence, uniqueness and Lipschitz dependence on initial data.
To the best of our knowledge, these are the first results on the Cauchy theory of two-dimensional Muskat problem with an elastic interface.
In particular, these are large-data well-posedness results that require only $\f12$ more derivative above the critical space.
Under this regularity setting, the curvature of the free surface $\Sigma_t$ may be unbounded and need not be locally square integrable.

ii) The restriction preventing local well-posedness in all subcritical spaces for arbitrary data arises from the remainder estimate for the Dirichlet–Neumann operator \eqref{GfEtaf}, which is valid only for $\mathfrak{s} \geq \f12$.
However, for sufficiently small initial data, the Muskat problem can be treated as a semilinear parabolic equation.
By exploiting the parabolic spacetime estimates \eqref{SpaceTimeOne} and \eqref{SpaceTimeTwo}, one can establish the global well-posedness of Muskat problem with an elastic interface in all sub-critical spaces, as stated in Theorem \ref{t:global} below. 

iii) As a quasilinear parabolic equation, the Muskat problem with an elastic interface exhibits the instantaneous parabolic smoothing.
Suppose $\eta \in C([0,T]; H^s(\R))$, then $\eta(t, \cdot) \in C^\infty(\R)$ when $t\in [\delta, T]$, for any $\delta>0$.
The proof is standard and therefore omitted. 
See \cite[Section 9]{MR3415681} for an analogous argument in the gravity Muskat setting.

\end{remark}

We define the space $\widetilde{L}^q(I; H^s(\R))$ as the space of tempered distributions $u \in \mathcal{S}' (\R^2)/ \mathcal{P}(\R)$ such that
\begin{equation*}
 \|u\|_{\widetilde{L}^q(I; H^s(\R))} : = \|2^{sj} \|P_j u \|_{L^q(I; L^2(\R))} \|_{\ell^2(\Z)} < \infty.
\end{equation*}
Our second main result is on the global well-posedness of  two-dimensional Muskat problem with an elastic interface in all sub-critical Sobolev spaces for small initial data.
 \begin{theorem}\label{t:global}
Let $s>\f32$, and the boundaries $\Gamma^\pm$ are either empty or flat $(\underline{b}^{\pm}_x = 0)$. 
We consider either the one-phase  or the  two-phase Muskat problem in the stable regime $\rho^+\leq \rho^-$. 
Let $\eta_0\in   H^s(\R)$, be an initial datum. Then there exist positive constants $\delta$ and $C$, such that the following holds: if $\| \eta_0\|_{H^s}\le \delta$ then for any $T>0$, \eqref{OneMuskat} and \eqref{TwoMuskatOne}-\eqref{TwoMuskatTwo} admit a unique solution $\eta\in C([0, T]; H^s (\R))$ satisfying
\begin{equation*}
\| \eta\|_{\wt L^\infty([0, T]; H^s)}+\frac{\sigma}{\mu^-} \| \eta\|_{\wt L^1([0, T];    H^{s+5})}\le C\| \eta_0\|_{H^s}.
\end{equation*}
 Moreover, if $\eta^1$ and $\eta^2$ to the Muskat problem, then for all $T>0$,
\begin{equation*}
\| \eta^1-\eta^2\|_{\wt L^\infty([0, T]; H^s)}+\frac{\sigma}{\mu^-} \| \eta^1-\eta^2\|_{\wt L^1([0, T]; H^{s+5})}\le C\| \eta^1(0)-\eta^2(0)\|_{H^s}.
\end{equation*}
\end{theorem}

For two-phase Muskat problem, the assumption $\rho^+ \leq  \rho^-$ corresponds to the situation where the  denser fluid lies below the less dense fluid, so that it is physically stable.
While local well-posedness holds regardless of density ratios, global well-posedness for small data is established here specifically for the stable regime $\rho^+ \leq  \rho^-$.
It is not clear whether or not one can obtain the global well-posedness result in the unstable regime for small data.
In the classical gravity Muskat problem without an elastic interface, the regime $\rho^+< \rho^-$ is known to be ill-posed in Sobolev spaces due to the heavier fluid being situated above the lighter one, leading to the rapid growth of interface perturbations.
See C\'ordoba \& Gancedo \cite{MR2318314} and \cite{MR2993754}, C\'ordoba, G\'omez-Serrano \& Zlato\v{s} \cite{MR3393318,MR3619874} for detailed results. 
In addition, for the Muskat problem with surface tension,  Lazar \cite{Lazar} established the global well-posedness in the strictly stable case $\rho^+<\rho^-$.

We summarize the main results established in the paper in the following table.
\begin{table}[ht]
\centering
\renewcommand{\arraystretch}{1.5}
\begin{tabular}{|l|l|c|l|}
\hline
\textbf{Phase} & \textbf{Result type} & \textbf{Regularity ($H^s(\R)$)} & \textbf{Data constraints} \\ \hline
one/two-phase & local well-posedness & $s \ge 2$ & arbitrary large data \\ \hline
one-phase & global well-posedness & $s > \frac{3}{2}$ & small initial data \\ \hline
two-phase (stable) & global well-posedness & $s > \frac{3}{2}$ & small data; $\rho^+ \le \rho^-$ \\ \hline
\end{tabular}
\caption{Well-posedness results for the Muskat problem established in this paper.}
\label{tab:muskat_results}
\end{table}

The proof of our main results relies on several key technical ingredients. 
First, we reformulate the Muskat problem into a single equation for the interface $\eta$ by utilizing the Dirichlet-Neumann operators $G^\pm(\eta)$ associated with the fluid domains $\Omega^\pm_t$. 
To address the high-order nonlinearities, we perform a detailed paralinearization of the elastic term $\mathbf{E}(\eta)$, which allows us to treat the Muskat problem as a fifth-order quasilinear paradifferential parabolic equation. 
For the local well-posedness, we derive a priori energy estimates and contraction estimates within the paradifferential framework to establish existence, uniqueness, and Lipschitz stability for arbitrary initial data. 
Finally, to obtain global results for small data, we rewrite the equation in an integral form based on a fifth-order fractional heat kernel and exploit parabolic space-time estimates to apply a fixed-point argument in sub-critical Sobolev spaces.

The rest of this paper is organized as follows.
In Section \ref{s:Reduct}, we first paralinearize the elastic term. 
Then we rewrite one-phase and two-phase Muskat problems as paradifferential parabolic equation and obtain a priori energy estimates for these equations.
In Section \ref{s:Contraction}, we first prove contraction estimates for both one-phase and two-phase Muskat problems, then we combine energy estimates and contraction estimates to obtain local well-posedness of Muskat problems.
Section \ref{s:Global} is then devoted to the proof of global well-posedness for Muskat problems.
We rewrite the solution of Muskat problems as fixed-point of some integral equations, and use the fixed-point Lemma \ref{t:fixedpoint} to obtain the unique existence of the global solution.
Many necessary paradifferential and product estimates are recalled in Appendix \ref{s:Paradifferential}.
Finally, we give some results on the Dirichlet-Neumann operators, the fractional heat kernel and the existence of fixed points in Appendix \ref{s:DNOperator}.

\section{Reduction of the Muskat problem and a priori estimates} \label{s:Reduct}
In this section we reformulate the Muskat problem using paradifferential calculus and derive a priori estimates for the Muskat problem. 
Throughout this section, we assume that $\eta \in Z^s(T)$ for $s\geq2$, and the surface remains at a positive distance from the boundaries $\Gamma^\pm$.
\begin{equation} \label{BottomH}
 \inf_{t\in[0,T]}\text{dist}(\eta(t), \Gamma^\pm)> h>0.
\end{equation}

We begin by paralinearizing the elastic term defined in $\mathbf{E}(\eta)$ in \eqref{Elastic}.
\begin{lemma} 
Let $s\geq2$, and $0<\delta \leq \f12$.
The elastic term admits the decomposition 
\begin{equation} \label{ElasticPara}
  \mathbf{E}(\eta) = T_{\ell} \eta + R_E(\eta),   
\end{equation}
where $ T_{\ell} \eta$ is the principal paradifferential term, and symbol $\ell$ is given by 
\begin{equation} \label{EllDef}
\begin{aligned}
    &\ell(x, \xi) = (1+\eta_x^2)^{-\f52}\xi^4 -2i((1+\eta_x^2)^{-\f52})_x \xi^3 \\
    -& \left(((1+\eta_x^2)^{-\f52})_{xx} -5(\eta_{xx}\eta_x(1+\eta_x^2)^{-\f72})_{x} +\f52\eta_{xx}^2(1-6\eta_x^2)(1+\eta_x^2)^{-\f92}\right) \xi^2 \\
    +& i\left(\f52 (\eta_{xx}^2(1-6\eta_x^2)(1+\eta_x^2)^{-\f92})_x-5(\eta_{xx}\eta_x(1+\eta_x^2)^{-\f72})_{xx}\right)\xi.
\end{aligned}    
\end{equation}
The lower-order remainder term $R_E(\eta)$ satisfies the estimate
\begin{equation} \label{REEtaEst}
    \|R_E(\eta)\|_{H^{s-\f32+\delta}} \lesssim \|\eta\|_{H^s} \|\eta\|_{H^{s+\f52}}. 
\end{equation}
\end{lemma}

\begin{proof}
Using the paralinearization \eqref{Paralinearization}, we can write
\begin{equation*}
\dfrac{\eta_x}{\sqrt{1+\eta_x^2}} =  T_{(1+\eta_x^2)^{-\f32}} \eta_x + R_1, \quad \dfrac{\eta_x}{(1+\eta_x^2)^\f72} =  T_{(1-6\eta_x^2)(1+\eta_x^2)^{-\f92}} \eta_x + R_2
\end{equation*}
where the remainder terms $R_1$ and $R_2$ satisfy
\begin{equation*}
 \|R_1\|_{H^{s+\f12 + \delta}} + \|R_2\|_{H^{s+\f12 + \delta}}\lesssim_{\|\eta_x\|_{L^\infty}} \| \eta_x\|_{H^{s+\f32}} \|\eta_x\|_{C^\delta_*} \lesssim \|\eta\|_{H^s} \| \eta\|_{H^{s+\f52}}.
\end{equation*}
We then get that using \eqref{BalancedError},
\begin{align*}
 &\left(\dfrac{\eta_x}{\sqrt{1+\eta_x^2}}\right)_x = T_{(1+\eta_x^2)^{-\f32}} \eta_{xx} -3T_{\eta_{xx}\eta_x(1+\eta_x^2)^{-\f52}} \eta_{x} + R_3, \quad \|R_3\|_{H^{s+\f12 + \delta}}\lesssim \|\eta\|_{H^s} \| \eta\|_{H^{s+\f52}}, \\
 &\dfrac{1}{1+\eta_x^2}\left(\dfrac{\eta_x}{\sqrt{1+\eta_x^2}}\right)_x = T_{(1+\eta_x^2)^{-\f52}} \eta_{xx} -5T_{\eta_{xx}\eta_x(1+\eta_x^2)^{-\f72}} \eta_{x} + R_4, \quad \|R_4\|_{H^{s+\f12 + \delta}}\lesssim \|\eta\|_{H^s} \| \eta\|_{H^{s+\f52}}, \\
 &\f{\eta_x\eta_{xx}^2}{(1+\eta^2_x)^{\f72}} = 2T_{\eta_{xx}\eta_x(1+\eta_x^2)^{-\f72}}\eta_{xx} + T_{\eta_{xx}^2(1-6\eta_x^2)(1+\eta_x^2)^{-\f92}} \eta_x + R_5, \quad \|R_5\|_{H^{s-\f12 + \delta}}\lesssim \|\eta\|_{H^s} \| \eta\|_{H^{s+\f52}}.
\end{align*}
By taking the spatial derivative on these terms, for the remainder term $R$ that satisfies \eqref{REEtaEst},
\begin{align*}
    &\left( \dfrac{1}{1+\eta_x^2}\left(\dfrac{\eta_x}{\sqrt{1+\eta_x^2}}\right)_x\right)_{xx} = T_{(1+\eta_x^2)^{-\f52}} \p_x^4\eta + 2T_{((1+\eta_x^2)^{-\f52})_x} \p_x^3\eta + T_{((1+\eta_x^2)^{-\f52})_{xx}} \eta_{xx}\\
    -&5T_{\eta_{xx}\eta_x(1+\eta_x^2)^{-\f72}} \p_x^3\eta -10T_{(\eta_{xx}\eta_x(1+\eta_x^2)^{-\f72})_x} \eta_{xx} -5T_{(\eta_{xx}\eta_x(1+\eta_x^2)^{-\f72})_{xx}} \eta_x + R, \\
    &\f52\left(\f{\eta_x\eta_{xx}^2}{(1+\eta^2_x)^{\f72}}\right)_x = 5T_{\eta_{xx}\eta_x(1+\eta_x^2)^{-\f72}}\p_x^3\eta + 5T_{(\eta_{xx}\eta_x(1+\eta_x^2)^{-\f72})_x}\eta_{xx}\\
    +& \f52 T_{\eta_{xx}^2(1-6\eta_x^2)(1+\eta_x^2)^{-\f92}} \eta_{xx}  + \f52 T_{(\eta_{xx}^2(1-6\eta_x^2)(1+\eta_x^2)^{-\f92})_x} \eta_{x} + R.
\end{align*}
Adding these two terms together, we get the paralinearization of the elastic term \eqref{ElasticPara}.
\end{proof}

We emphasize that the precise explicit expression of the symbol $\ell$ is not essential to the subsequent analysis.
What is crucial is its elliptic structure.
In particular, for $u\in H^{s+4}$, $s> \f32$, the symbol $\ell$ satisfies
\begin{equation} \label{ellEst}
    \|T_{\ell} u - \p_x^4 u\|_{H^s} \lesssim \|\eta_x\|_{C^\epsilon_*} \|u \|_{H^{s+4}} \lesssim \|\eta\|_{H^s} \| \eta\|_{H^{s+4}}.
\end{equation}

\subsection{The one-phase case}
The one-phase Muskat problem \eqref{OneMuskat} is reformulated into a paradifferential form to facilitate energy estimates.
\begin{proposition}
For $\delta \in (0, \f12]$, the one-phase Muskat problem \eqref{OneMuskat} can be rewritten as 
\begin{equation} \label{ParaOneMuskat}
 \p_t \eta = -\frac{\sigma}{\mu^{-}} |D|T_\ell \eta + K, 
\end{equation}
where the remainder term $K$ is controlled by the initial data and gravitational terms as follows:
\begin{equation*}
    \|K\|_{H^{s-\f52+\delta}} \lesssim_{\|\eta\|_{H^s}} \frac{\sigma}{\mu^{-}}  \|\eta\|_{H^{s+\f52}} + \frac{\rho^- g}{\mu^-}\| \eta\|_{H^{s-\f32+\delta}}.
\end{equation*}
\end{proposition}

\begin{proof}
Recall the one-phase Muskat problem reads
\begin{equation*}
  \p_t \eta = -\frac{\sigma}{\mu^-}G^-(\eta) \mathbf{E}(\eta) -\frac{\rho^-g}{\mu^-} G^-(\eta)\eta.
\end{equation*}
For the term $G^-(\eta)\eta$, the estimate \eqref{GfEtaf} with $\mathfrak{s} = s-\f32+\delta$ yields
\begin{equation*}
   \| G^-(\eta)\eta\|_{H^{s-\f52+\delta}} \lesssim_{\|\eta\|_{H^s}} \|\eta\|_{H^{s-\f32+\delta}}.
\end{equation*}
Regarding the term $G^-(\eta) \mathbf{E}(\eta)$, we use \eqref{GfDRf}, and \eqref{GfEtaf} to write
\begin{equation*}
 G^-(\eta) \mathbf{E}(\eta) = |D|\mathbf{E}(\eta) + R^{-}(\eta)\mathbf{E}(\eta),
\end{equation*}
where using \eqref{GfDRf},
\begin{equation*}
 \| R^{-}(\eta)\mathbf{E}(\eta) \|_{H^{s-\f52+\delta}} \lesssim_{\|\eta\|_{H^s}} \|\mathbf{E}(\eta) \|_{H^{s-\f32}} \lesssim_{\|\eta\|_{H^s}} \|\eta\|_{H^{s+\f52}},
\end{equation*}
so that this term can be put into $K$.
For the $|D|\mathbf{E}(\eta)$ term, we use \eqref{ElasticPara} to estimate the difference
\begin{equation*}
\||D|\mathbf{E}(\eta) - |D|T_{\ell}\eta \|_{H^{s-\f52+\delta}} \lesssim  \|\mathbf{E}(\eta) - T_{\ell}\eta \|_{H^{s-\f32+\delta}}\lesssim_{\|\eta\|_{H^s}} \|\eta\|_{H^{s+\f52}}.    
\end{equation*}
By putting the perturbative terms into $K$, we obtain the estimate \eqref{ParaOneMuskat}.
\end{proof}

At leading order, the symbol $|\xi|\ell(x,\xi) \approx (1+\eta_x^2)^{-\f52}|\xi|^5$, which is elliptic of order $5$.
Consequently, the one-phase Muskat problem \eqref{ParaOneMuskat} is a fifth-order quasilinear parabolic equation.

We have the following a priori estimate for the one-phase Muskat equation.
\begin{proposition} \label{t:OnePhaseEst}
Let $s\geq2$, assume that $\eta \in Z^s(T)$ is a solution to \eqref{ParaOneMuskat} that satisfies \eqref{BottomH}.
Then there exists a function $\mathcal{F}: \R^+ \rightarrow \R^+$ depending only on $(h,s,\frac{\rho^- g}{\mu^-},\frac{\sigma}{\mu^-})$ such that
\begin{equation} \label{OneAprioriEst}
  \| \eta\|_{Z^s(T)} \lesssim \mathcal{F}\big(\|\eta(0,\cdot)\|_{H^s} + T^\f12\mathcal{F}( \| \eta\|_{Z^s(T)})\big).
\end{equation}
\end{proposition}

\begin{proof}
We apply $\langle T_{(1+\eta_x^2)^{-\frac{s}{2}}}D\rangle^s$ to the equation \eqref{ParaOneMuskat}, and denote $\eta_s := \langle T_{(1+\eta_x^2)^{-\frac{s}{2}}}D\rangle^s \eta$.
Then $\eta_s$ solves the equation
\begin{equation*} 
     \p_t \eta_s = -\frac{\sigma}{\mu^{-}} |D|T_\ell \eta_s -\frac{\sigma}{\mu^{-}}\big[\langle T_{(1+\eta_x^2)^{-\frac{s}{2}}}D \rangle^s, |D|T_\ell \big] \eta_s +  \langle T_{(1+\eta_x^2)^{-\frac{s}{2}}}D\rangle^sK.
\end{equation*}
Taking the time derivative of the $\|\eta_s\|_{L^2}^2$, we get 
\begin{align*}
 \f12\frac{d}{dt} \|\eta_s\|_{L^2}^2 =& -\frac{\sigma}{\mu^{-}} (|D|T_\ell \eta_s, \eta_s)_{L^2 \times L^2} - \frac{\sigma}{\mu^{-}} ([\langle T_{(1+\eta_x^2)^{-\frac{s}{2}}}D \rangle^s, |D|T_\ell] \eta_s, \eta_s)_{L^2 \times L^2} \\
 &+ (\langle T_{(1+\eta_x^2)^{-\frac{s}{2}}}D\rangle^sK, \eta_s)_{L^2 \times L^2}.
\end{align*}
From the above proposition, for the last term involving $K$,
\begin{align*}
&\left|(\langle T_{(1+\eta_x^2)^{-\frac{s}{2}}}D\rangle^sK, \eta_s)_{L^2 \times L^2}\right|
\leq\|\langle T_{(1+\eta_x^2)^{-\frac{s}{2}}}D\rangle^sK, \eta_s\|_{H^{-\f52+\delta}}\|\eta_s\|_{H^{\f52-\delta}}\\
&\lesssim_{\|\eta\|_{H^s}}\left(\frac{\sigma}{\mu^{-}}   + \frac{\rho^- g}{\mu^-}\right)\|\eta\|_{H^{s+\f52}}\| \eta\|_{H^{s+\f52-\delta}}.
\end{align*}
Since at the leading order,  we use the symbolic calculus \eqref{CompositionPara}, the commutator satisfies
\begin{equation*}
 \big\|[\langle T_{(1+\eta_x^2)^{-\frac{s}{2}}}D \rangle^s, |D|T_\ell] \eta_s \big\|_{H^{-\f52+\delta}} \lesssim_{\|\eta\|_{H^s}} \|\eta\|_{H^{s+\f52}},
\end{equation*}
so that the commutator term has the estimate
\begin{equation*}
\big| ([\langle T_{(1+\eta_x^2)^{-\frac{s}{2}}}D \rangle^s, |D|T_\ell] \eta_s, \eta_s)_{L^2 \times L^2}\big| \lesssim_{\|\eta\|_{H^s}}\|\eta\|_{H^{s+\f52}}\| \eta\|_{H^{s+\f52-\delta}}.
\end{equation*}
For the first term in the time derivative of the $\|\eta_s\|_{L^2}^2$, we split the symbol as the sum of the leading order term, and the lower-order term 
\begin{equation*}
    |D|T_\ell =  T_{m^{[5]}} + T_{m^{[\leq 4]}}, \quad m^{[5]} = (1+\eta_x^2)^{-\f52}|\xi|^5,
\end{equation*}
where the symbol $m^{[\leq 4]}$ is the lower-order part of $|D|T_\ell$ of order at most $4$, and $m^{[5]}$ is strictly elliptic as long as $\eta \in \dot{W}^{1,\infty}$.
For the estimate that involves $m^{[\leq 4]}$, we get 
\begin{equation*}
 \frac{\sigma}{\mu^{-}} |(T_{m^{[\leq 4]}} \eta_s, \eta_s)_{L^2 \times L^2}| \leq \frac{\sigma}{\mu^{-}} \| T_{m^{[\leq 4]}} \eta_s\|_{H^{-\f52}} \| \eta_s\|_{H^{\f52}}  \lesssim_{\|\eta\|_{H^s}} \frac{\sigma}{\mu^{-}}\|\eta\|_{H^{s+\f52}}\| \eta\|_{H^{s+\f52-\delta}}.     
\end{equation*}
For the term that involves $m^{[5]}$, we write 
\begin{align*}
 (T_{m^{[5]}} \eta_s, \eta_s)_{L^2 \times L^2} &= \big(T_{\sqrt{m^{[5]}}} \eta_s, T_{\sqrt{m^{[5]}}}\eta_s\big)_{L^2 \times L^2} + \big(T_{\sqrt{m^{[5]}}} \eta_s, \big((T_{\sqrt{m^{[5]}}})^{*} -T_{\sqrt{m^{[5]}}}\big)\eta_s\big)_{L^2 \times L^2} \\
 &+ \big((T_{m^{[5]}}-T_{\sqrt{m^{[5]}}} T_{\sqrt{m^{[5]}}}) \eta_s, \eta_s\big)_{L^2 \times L^2}.
\end{align*}
For the first term in $(T_{m^{[5]}} \eta_s, \eta_s)_{L^2 \times L^2}$, following the estimate $(3.18)$ in \cite{MR4131404}, by uisng \eqref{TABound} and \eqref{CompositionPara}, we have
\begin{align*}
\|\eta_s\|_{H^{\f52}}&\leq
\|T_{\sqrt{m^{[5]}}^{-1}}T_{\sqrt{m^{[5]}}}\eta_s\|_{H^{\f52}}
+\|(Id-T_{\sqrt{m^{[5]}}^{-1}}T_{\sqrt{m^{[5]}}})\eta_s\|_{H^{\f52}}\\
&\lesssim_{\|\eta\|_{H^s}}\|T_{\sqrt{m^{[5]}}}\eta_s\|_{L^2}+\|\eta_s\|_{H^{\f52-\delta}},
\end{align*}
which gives that
\begin{equation*}
  \|T_{\sqrt{m^{[5]}}} \eta_s\|^2_{L^2} \geq \frac{1}{\mathcal{F}(\|\eta\|_{H^s})} \|\eta\|^2_{H^{s+\f52}} - \mathcal{F}(\|\eta\|_{H^s})\|\eta\|_{H^{s+\f52}}\| \eta\|_{H^{s+\f52-\delta}},
\end{equation*}
function $\mathcal{F}: \R^+ \rightarrow \R^+$.
Using the symbolic calculus \eqref{CompositionPara} and \eqref{AdjointBound}, 
\begin{equation*}
 \|\big((T_{\sqrt{m^{[5]}}})^{*} -T_{\sqrt{m^{[5]}}}\big)\eta_s \|_{H^\delta} + \|(T_{m^{[5]}}-T_{\sqrt{m^{[5]}}} T_{\sqrt{m^{[5]}}}) \eta_s\|_{H^{-\f52 + \delta}} \lesssim_{\|\eta\|_{H^s}}\|\eta\|_{H^{s+\f52}}.
\end{equation*}
we bound the second and the third terms in $(T_{m^{[5]}} \eta_s, \eta_s)_{L^2 \times L^2}$,
\begin{align*}
    &\big(T_{\sqrt{m^{[5]}}} \eta_s, \big((T_{\sqrt{m^{[5]}}})^{*} -T_{\sqrt{m^{[5]}}}\big)\eta_s\big)_{L^2 \times L^2} + \big((T_{m^{[5]}}-T_{\sqrt{m^{[5]}}} T_{\sqrt{m^{[5]}}}) \eta_s, \eta_s\big)_{L^2 \times L^2}\\
    &\lesssim_{\|\eta\|_{H^s}}\|\eta\|_{H^{s+\f52}}\| \eta\|_{H^{s+\f52-\delta}}.
\end{align*}
Combining the above estimates 
\begin{equation*}
-(|D|T_\ell \eta_s, \eta_s)_{L^2 \times L^2} \leq - \frac{1}{\mathcal{F}(\|\eta\|_{H^s})} \|\eta\|^2_{H^{s+\f52}} +\mathcal{F}(\|\eta\|_{H^s})\|\eta\|_{H^{s+\f52}}\| \eta\|_{H^{s+\f52-\delta}}. 
\end{equation*}
As a consequence,
\begin{equation*}
 \f12\frac{d}{dt} \|\eta_s\|_{L^2}^2 \leq -\frac{\sigma}{\mu^{-}}\frac{1}{\mathcal{F}(\|\eta\|_{H^s})} \|\eta\|^2_{H^{s+\f52}} + \Big(\frac{\sigma}{\mu^{-}} + \frac{\rho^- g}{\mu^-}  \Big) \mathcal{F}(\|\eta\|_{H^s})\|\eta\|_{H^{s+\f52}}\| \eta\|_{H^{s+\f52-\delta}},
\end{equation*}
for some function $\mathcal{F}$ that depends only on $h$ and $s$.
For the factor $\| \eta\|_{H^{s+\f52-\delta}}$, we use the interpolation
\begin{equation*}
 \| \eta\|_{H^{s+\f52-\delta}} \lesssim \| \eta\|_{H^s}^{1-\theta}\| \eta\|_{H^{s+\f52}}^{\theta},
\end{equation*}
for some constant $\theta \in (0,1)$ depending on $\delta$ and $s$.
Using Young's inequality,  we then get
\begin{align*}
\f12\f{d}{dt}\|\eta_s\|^2_{L^2}
\leq-\f{1}{\mathcal{F}(\|\eta\|_{H^s})}\|\eta\|^2_{H^{s+\f52}}+\mathcal{F}(\|\eta\|_{H^s})\|\eta\|^2_{H^{s}},
\end{align*}
which implies that for $t\in(0,T]$
\begin{align*}
\f{d}{dt}\|\eta_s\|^2_{L^2}
\leq-\f{1}{\mathcal{F}\big(\|\eta\|_{Z^s(T)} \big)}\|\eta\|^2_{H^{s+\f52}}+\mathcal{F}\big(\|\eta\|_{Z^s(T)}\big)\|\eta\|^2_{H^{s}}.
\end{align*}
Hence, taking the time integral, for $t\in(0,T]$
\begin{align*}
\|\eta_s(t,\cdot)\|^2_{L^2}+\f{1}{\mathcal{F}\big(\|\eta\|_{Z^s(T)} \big)}\int_0^t\|\eta(s,\cdot)\|^2_{H^{s+\f52}}ds\leq \|\eta_s(0,\cdot)\|^2_{L^2}+T\mathcal{F}\big(\|\eta\|_{Z^s(T)}\big)\|\eta\|^2_{L^\infty_t H^{s}_x}.
\end{align*}
Since we have the norm equivalence
\begin{align*}
\f{1}{\mathcal{F}\big(\|\eta\|_{Z^s(T)} \big)}\|\eta(t,\cdot)\|_{H^s}\leq\|\eta_s(t,\cdot)\|_{L^2}\leq\mathcal{F}\big(\|\eta\|_{Z^s(T)}\big)\|\eta(t,\cdot)\|_{H^s},
\end{align*}
we then obtain that for $t\in(0,T]$,
\begin{align*}
\|\eta(t,\cdot)\|^2_{H^s}+\int_0^t\|\eta(s,\cdot)\|^2_{H^{s+\f52}}ds\leq \mathcal{F}\left(\|\eta\|_{Z^s(T)}\right)\left(\|\eta(0, \cdot)\|^2_{H^s}+T\mathcal{F}\big(\|\eta\|_{Z^s(T)}\big)\|\eta\|^2_{Z^s(T)}\right),
\end{align*}
which can be written as 
\begin{align*}
\|\eta\|_{Z^s(T)}^2\leq \mathcal{F}\left(\|\eta\|_{Z^s(T)}\right)\left(\|\eta(0, \cdot)\|_{H^s}^2+T\mathcal{F}\big(\|\eta\|_{Z^s(T)}\big)\|\eta\|^2_{Z^s(T)}\right).
\end{align*}
This leads to the a priori estimate \eqref{OneAprioriEst}.
\end{proof}


We also need an a priori estimate for the distance of the surface $\Sigma_t$ and the bottom $\Gamma^{-}$.
\begin{lemma} \label{t:DistBound}
Let $s\geq 2$, and $\eta \in Z^s(T)$ is a solution to \eqref{ParaOneMuskat} that satisfies \eqref{BottomH}.
Then there exist $\theta \in (0,1)$ and a function $\mathcal{F}: \R^+ \rightarrow \R^+$ depending only on $(h,s,\frac{\rho^- g}{\mu^-},\frac{\sigma}{\mu^-})$ such that
\begin{equation} \label{DistBound}
    \inf_{t\in[0,T]} \text{dist}(\eta(t), \Gamma^{-}) \geq \text{dist}(\eta(0), \Gamma^{-}) - T^\theta \mathcal{F}(\| \eta\|_{Z^s(T)}).
\end{equation}
In particular, for $T>0$ sufficiently small, the free surface $\Gamma_t$ remains uniformly separated from the rigid boundary $\Gamma^-$.
\end{lemma}

\begin{proof}
Using the equation \eqref{OneMuskat}, the estimate \eqref{GfEtaf} and the fact that $s+\f52>\f{9}{2}$, we estimate
 \begin{align*}
 &\|\eta(t) - \eta(0)\|_{H^{-\f12}}\leq \int_0^t \frac{\sigma}{\mu^{-}} \|G^{-}(\eta)H(\eta)(\tau)\|_{H^{-\f12}} + \frac{\rho^{-}g}{\mu^{-}} \|G^{-}(\eta)\eta(\tau)\|_{H^{-\f12}} d\tau \\
 \lesssim&  \int_0^t \mathcal{F}(\|\eta\|_{H^s}) \| \eta\|_{H^\f92}d\tau \leq t^\f12 \mathcal{F}\big(\|\eta\|_{L^\infty([0,t];H^s)}\big)\|\eta\|_{L^2([0,t];H^{s+\f52})}.
 \end{align*}
 For any constant $\nu \in (2, s)$, we use interpolation of Sobolev spaces:
\begin{equation*}
 \|\eta(t) - \eta(0)\|_{H^\nu} \leq \|\eta(t) - \eta(0)\|_{H^{-\f12}}^\theta \|\eta(t) - \eta(0)\|_{H^{s}}^{1-\theta} \leq t^{\frac{\theta}{2}} \mathcal{F}\big(\|\eta\|_{Z^s(t)}\big),
\end{equation*}
for some constant $\theta \in (0,1)$ that depends on $s$ and $\nu$.
Using the Sobolev embedding $H^\nu(\R) \subset L^\infty(\R)$, we obtain the estimate \eqref{DistBound}.
\end{proof}

\subsection{The two-phase case}
For two-phase Muskat problem, we first recall a well-posedness result of \eqref{TwoMuskatTwo} from Proposition $4.1$ in \cite{MR2812947}.
\begin{lemma} [\hspace{1sp}\cite{MR2812947}] \label{t:fpmHHalfEst}
Let $\eta \in W^{1,\infty}(\R)\cap H^\f12(\R)$ satisfy dist$(\eta, \Gamma^\pm) >h>0$.
Then there exists a unique variational solution $f^\pm \in \tilde{H}^\f12_\pm (\R)$ to the system \eqref{TwoMuskatTwo}.
Moreover, $f^\pm$ satisfy
\begin{equation*} 
 \|f^\pm\|_{\tilde{H}^\f12_\pm} \leq C(1+\|\eta\|_{W^{1,\infty}})^2 \| \sigma \mathbf{E}(\eta) + g(\rho^- - \rho^+)\eta\|_{H^\f12},
\end{equation*}
where the constant $C$ depends only on $(h, \mu^\pm)$.
\end{lemma}
Using the estimate \eqref{ellEst}, the above estimate is reduced to
\begin{equation} \label{fpmEst}
 \|f^\pm\|_{\tilde{H}^\f12_\pm} \lesssim_{\|\eta\|_{W^{1,\infty}}}  \sigma \|\eta\|_{H^{\f92}} + g(\rho^- - \rho^+)\|\eta\|_{H^\f12}.
\end{equation}

For the profile $\eta$ with higher Sobolev regularity, we change the estimate for the capillary term to the elastic term in Proposition $4.2$ in \cite{MR2812947}, and obtain the following result.
\begin{lemma} [\hspace{1sp}\cite{MR2812947}]
Let $f^\pm$ be the solution as given by Lemma \ref{t:fpmHHalfEst}.
If $\eta \in H^{s+\f52}$ with $s\geq 2$ then $f^\pm \in \tilde{H}^{s-\f12}_{\pm}(\R)$, and 
\begin{equation} \label{fpmEstR}
 \|f^\pm\|_{\tilde{H}^r_\pm} \lesssim_{\|\eta\|_{H^s}}  \sigma \|\eta\|_{H^{r+4}} + g(\rho^- - \rho^+)\|\eta\|_{H^r},
\end{equation}
for all $r\in [\f12, s-\f12]$, where the function $\mathcal{F}$ depends only on $(h,s,r, \mu^\pm)$.
\end{lemma}

Then we rewrite the two-phase Muskat problem \eqref{TwoMuskatOne}-\eqref{TwoMuskatTwo} in the paradifferential form.
\begin{proposition}
 For $\delta \in (0,\f12]$, the two-phase Muskat problem can be written as
 \begin{equation} \label{ParaTwoMuskat}
    \p_t \eta = -\frac{\sigma}{\mu^+ +\mu^{-}} |D|T_\ell \eta + K,
 \end{equation}
 where the remainder term $K$ satisfies the estimate
\begin{equation*}
    \|K\|_{H^{s-\f52+\delta}} \lesssim_{\|\eta\|_{H^s}} \sigma  \|\eta\|_{H^{s+\f52}} + (\rho^- - \rho^+)g\| \eta\|_{H^{s-\f32+\delta}}.
\end{equation*}
\end{proposition}

\begin{proof}
According to the estimate \eqref{GfDRf}, for $\delta \in (0,\f12]$, and $\mathfrak{s} \in [\f12, s-\delta]$, there exists some function $\mathcal{F}:\R^+ \rightarrow \R^+$ depending only on $(s,\mathfrak{s}, \delta, h)$ such that
\begin{equation*} 
G^\pm(\eta)f^\pm = \mp|D|f^\pm + R^\pm(\eta)f^\pm, \quad \| R^\pm(\eta)f^\pm\|_{H^{\mathfrak{s}- 1+\delta}} \leq \mathcal{F}(\|\eta\|_{H^s}) \|f\|_{\tilde{H}^\mathfrak{s}_-}.
\end{equation*}
 We then write the system \eqref{TwoMuskatTwo} as 
\begin{align*}
 |D|f^- &= \frac{\mu^-}{\mu^+ + \mu^-} |D|(\sigma \mathbf{E}(\eta)+ (\rho^- - \rho^+)g\eta) + \frac{\mu^-}{\mu^+ + \mu^-}R^+(\eta)f^+ - \frac{\mu^+}{\mu^+ + \mu^-}R^-(\eta)f^- \\
 &= \frac{\sigma\mu^-}{\mu^+ + \mu^-} |D|T_\ell \eta + \frac{\sigma\mu^-}{\mu^+ + \mu^-} |D|(\mathbf{E}(\eta)-T_\ell \eta) + \frac{\mu^-  (\rho^- - \rho^+)g}{\mu^+ + \mu^-} |D|\eta\\
 &+ \frac{\mu^-}{\mu^+ + \mu^-}R^+(\eta)f^+ - \frac{\mu^+}{\mu^+ + \mu^-}R^-(\eta)f^-.
\end{align*}
We then estimate using \eqref{ellEst} and \eqref{fpmEstR},
\begin{align*}
 &\| |D|(\mathbf{E}(\eta)-T_\ell \eta)\|_{H^{s-\f52 + \delta}} \lesssim_{\|\eta\|_{H^s}} \|\eta\|_{H^{s+\f32}}, \\
 & \||D|\eta \|_{H^{s-\f52 + \delta}} \lesssim_{\|\eta\|_{H^s}} \|\eta\|_{H^{s-\f32+\delta}}, \\
 & \|R^\pm(\eta) f^\pm\|_{H^{s-\f52+\delta}}\lesssim_{\|\eta\|_{H^s}}\|f^\pm\|_{\tilde{H}_\pm^{s-\f32}} \lesssim_{\|\eta\|_{H^s}}  \sigma \|\eta\|_{H^{s+\f52}} + g(\rho^- - \rho^+)\|\eta\|_{H^{s-\f32}}.
\end{align*}
Hence, we obtain that
\begin{equation*}
|D|f^- = \frac{\sigma\mu^-}{\mu^+ + \mu^-} |D|T_\ell \eta  + K.
\end{equation*}
Using the equation \eqref{TwoMuskatOne}, we get
\begin{equation*}
\p_t \eta = -\frac{1}{\mu^-}|D|f^- -\frac{1}{\mu^-}R^-(\eta)f^- = -\frac{\sigma}{\mu^+ +\mu^{-}} |D|T_\ell \eta + K,
\end{equation*}
which gives the reformulation of the two-phase Muskat problem \eqref{ParaTwoMuskat}.
\end{proof}

Similar to the one-phase problem, as long as $\eta \in \dot{W}^{1,\infty}$, the two-phase Muskat problem \eqref{ParaTwoMuskat} is a fifth-order quasilinear parabolic equation.
Since the paradifferential form of the two-phase Muskat problem \eqref{ParaTwoMuskat} is similar to the paradifferential form of the one-phase Muskat problem \eqref{ParaOneMuskat}, we follow the computation in Proposition \ref{t:OnePhaseEst} and \ref{t:DistBound}, and obtain the a priori estimate for the two-phase Muskat problem.
\begin{proposition} \label{t:TwoPhaseEst}
Let $s\geq 2$, and suppose that $\eta$ solves the paradifferntial two-phase Muskat problem \eqref{ParaTwoMuskat} on $[0,T]$, and satisfies the separation condition \eqref{BottomH}.
Then there exist $\theta\in (0,1)$ depending only on $s$, and a function $\mathcal{F}: \R^+ \rightarrow \R^+$ depending only on $(h,s, \sigma, \mu^\pm, (\rho^- -\rho^+)g)$ such that
\begin{equation*} 
  \| \eta\|_{Z^s(T)} \lesssim \mathcal{F}\big(\|\eta(0,\cdot)\|_{H^s} + T^\f12\mathcal{F}( \| \eta\|_{Z^s(T)})\big).
\end{equation*}
Moreover, 
\begin{equation*} 
    \inf_{t\in[0,T]} \text{dist}(\eta(t), \Gamma^{-}) \geq \text{dist}(\eta(0), \Gamma^{-}) - T^\theta \mathcal{F}\big(\| \eta\|_{Z^s(T)}\big).
\end{equation*}
\end{proposition}

\section{Contraction estimates and proof of local well-posedness} \label{s:Contraction}
This section is dedicated to proving the contraction estimates for both the one-phase and two-phase Muskat problems.
These results are then combined with the a priori estimates from Section \ref{s:Reduct} to establish Theorem \ref{t:MainOne}.

We begin with a contraction estimate for the remainder term in the paralinearization of the elastic term $R_E(\eta)$ in \eqref{ElasticPara}.

\begin{lemma} \label{t:REEst}
For $\delta\in (0,s- \f32)$, and $\delta \leq 1$, $R_E(\eta)$ satisfies the contraction estimate
\begin{equation*}
 \|R_E(\eta_1) - R_E(\eta_2) \|_{H^{s-\f32}} \lesssim_{\|(\eta_1, \eta_2)\|_{H^s \times H^s}}  \Big(\|\eta_1\|_{H^{s+\f52}}+ \|\eta_2\|_{H^{s+\f52}} \Big) \| \eta_1 -\eta_2\|_{H^s}.
\end{equation*}
\end{lemma}
\begin{proof}
We denote the G\^{a}teaux derivative $d_uF(u)$ of a function $F$ at $u$ in the direction $\dot{u}$ as follows:
\begin{align*}
d_uF(u)\dot{u}=\lim_{\e\to0}\f{1}{\e}(F(u+\e\dot{u})-F(u)).
\end{align*}
To obtain this lemma, we just need prove
\begin{align} \label{dEtaREEta}
\|d_\eta R_{E}(\eta)\dot{\eta}\|_{H^{s-\f32}}\lesssim_{\|\eta\|_{H^s}}  \|\eta\|_{H^{s+\f52}} \| \dot{\eta}\|_{H^s}.
\end{align}
One can check that
\begin{align*}
&d_{\eta}\left(\f{\eta_x}{\sqrt{1+\eta_x^2}}\right)\dot{\eta}=\f{\dot{\eta}_x}{(1+\eta_x^2)^{\f32}},\quad
d_{\eta}\left(\f{1}{1+\eta_x^2}\right)\dot{\eta}=-\f{2\eta_x\dot{\eta}_x}{(1+\eta_x^2)^2},\\
&d_{\eta}\left(\f{\eta_x}{(1+\eta_x^2)^{\f72}}\right)\dot{\eta}=\f{1-6\eta_x^2}{(1+\eta_x^2)^{\f92}}\dot{\eta}_x.
\end{align*}
Hence, using the chain rule for the G\^{a}teaux derivative,
\begin{align*}
d_{\eta}\left[\f{1}{1+\eta_x^2}\left(\dfrac{\eta_x}{\sqrt{1+\eta_x^2}}\right)_x \right]\dot{\eta}
=\f{\dot{\eta}_{xx}}{(1+\eta^2_x)^{\f52}}-\f{5\eta_x\eta_{xx}\dot{\eta}_x}{(1+\eta_x^2)^{\f72}},\\
d_{\eta}\left(\f{\eta_x\eta_{xx}^2}{(1+\eta^2_x)^{\f72}}\right)\dot{\eta}
=\f{(1-6\eta_x^2)\eta^2_{xx}}{(1+\eta_x^2)^{\f92}}\dot{\eta}_x
+\f{2\eta_x\eta_{xx}}{(1+\eta_x^2)^{\f72}}\dot{\eta}_{xx}.
\end{align*}
Therefore, we have
\begin{align*}
&d_{\eta}\left(\f{1}{1+\eta_x^2}\left(\dfrac{\eta_x}{\sqrt{1+\eta_x^2}}\right)_x\right)_{xx}\dot{\eta}=\f{\p_x^4\dot{\eta}}{(1+\eta_x)^{\f52}}-\f{10\eta_x\eta_{xx}\p^3_x\dot{\eta}}{(1+\eta_x^2)^{\f72}}\\
+&((1+\eta^2_x)^{-\f52})_{xx}\dot{\eta}_{xx}-\f{5\eta_x\eta_{xx}\dot{\eta}_{xxx}}{(1+\eta_x^2)^{\f72}}-\left(\f{10\eta_x\eta_{xx}}{(1+\eta_x^2)^{\f72}}\right)_x\dot{\eta}_{xx}-\left(\f{5\eta_x\eta_{xx}}{(1+\eta_x^2)^{\f72}}\right)_{xx}\dot{\eta}_{x},
\end{align*}
and
\begin{align*}
\f52d_{\eta}\left(\f{\eta_x\eta_{xx}^2}{(1+\eta^2_x)^{\f72}}\right)_x\dot{\eta}
=&\f{5\eta_x\eta_{xx}}{(1+\eta_x^2)^{\f72}}\p^3_{x}\dot{\eta}+\left(\f{5\eta_x\eta_{xx}}{(1+\eta_x^2)^{\f72}}\right)_x\dot{\eta}_{xx}\\
+&\f{5(1-6\eta_x^2)\eta^2_{xx}}{2(1+\eta_x^2)^{\f92}}\dot{\eta}_{xx}+\left(\f{5(1-6\eta_x^2)\eta^2_{xx}}{2(1+\eta_x^2)^{\f92}}\right)_x\dot{\eta}_{x}.
\end{align*}
Thus, we obtain that
\begin{align*}
&d_{\eta}\mathbf{E}(\eta)\dot{\eta}=
\f{\p_x^4\dot{\eta}}{(1+\eta^2_x)^{\f52}}+2\big((1+\eta^2_x)^{-\f52}\big)_x\p^3_x\dot{\eta}\\
&+\left\{((1+\eta^2_x)^{-\f52})_{xx}-\left(\f{5\eta_x\eta_{xx}}{(1+\eta_x^2)^{\f72}}\right)_x+\f{5(1-6\eta_x^2)\eta^2_{xx}}{2(1+\eta_x^2)^{\f92}}\right\}\dot{\eta}_{xx}\\
&-\left[\left(\f{5\eta_x\eta_{xx}}{(1+\eta_x^2)^{\f72}}\right)_{xx}-\left(\f{5(1-6\eta_x^2)\eta^2_{xx}}{2(1+\eta_x^2)^{\f92}}\right)_x\right]\dot{\eta}_{x}.
\end{align*}
Recall that
\begin{align*}
   T_{\ell}\dot{\eta} =&T_{(1+\eta_x^2)^{-\f52}} \p_x^4\dot{\eta} + 2T_{((1+\eta_x^2)^{-\f52})_x} \p_x^3\dot{\eta}+ T_{((1+\eta_x^2)^{-\f52})_{xx}-5(\eta_{xx}\eta_x(1+\eta_x^2)^{-\f72})_x+\f52\eta_{xx}^2(1-6\eta_x^2)(1+\eta_x^2)^{-\f92}} \dot{\eta}_{xx}\\
+&T_{-5\eta_{xx}\eta_x(1+\eta_x^2)^{-\f72}+\f52(\eta_{xx}^2(1-6\eta_x^2)(1+\eta_x^2)^{-\f92})_x} \dot{\eta}_x.
\end{align*}
By using \eqref{BalancedError}, we then get
\begin{align*}
\|d_{\eta}\mathbf{E}(\eta)\dot{\eta}-T_{\ell}\dot{\eta}\|_{H^{s-\f32}}\lesssim_{\|\eta\|_{H^s}}   \|\eta\|_{H^{s+\f52}} \| \dot{\eta}\|_{H^s}.
\end{align*}
Since 
\begin{align*}
d_{\eta}R_E(\eta)\dot{\eta}=d_{\eta}\mathbf{E}(\eta)\dot{\eta}-T_{\ell}\dot{\eta}-T_{d_{\eta}\ell\dot{\eta}}\eta,
\end{align*}
it remains to prove
\begin{equation} \label{tdEtaEllEtaEta}
\|T_{d_{\eta}\ell\dot{\eta}}\eta\|_{H^{s-\f32}}\lesssim_{\|\eta\|_{H^s}} \|\eta\|_{H^{s+\f52}} \| \dot{\eta}\|_{H^s}.
\end{equation}
This can be deduced from \eqref{TABound}. Thus, we prove \eqref{dEtaREEta}, which leads to the proof of the lemma.
\end{proof}
\begin{lemma}
For $0<\delta \leq \f12$ and $s\geq2$, $\mathbf{E}(\eta)$ satisfies the contraction estimate
\begin{equation} \label{ElasticDiff}
\begin{aligned}
 \|\mathbf{E}(\eta_1) - \mathbf{E}(\eta_2) \|_{H^{s-\f32-\delta}} &\lesssim_{\|(\eta_1,\eta_2)\|_{H^s}}\|\eta_1 - \eta_2\|_{H^{s+\f52-\delta}}\\
 +&(\|\eta_1\|_{H^{s+\f52-\delta}}+\|\eta_2\|_{H^{s+\f52-\delta}})\|\eta_1 - \eta_2\|_{H^{s}}.
 \end{aligned}
\end{equation}
\end{lemma}
\begin{proof}
Recall that
\begin{equation*}
d_{\eta}\mathbf{E}(\eta)\dot{\eta} = d_{\eta}R_E(\eta)\dot{\eta}+T_{\ell}\dot{\eta}+T_{d_{\eta}\ell\dot{\eta}}\eta.
\end{equation*}
Then we use \eqref{ellEst}, \eqref{dEtaREEta} and \eqref{tdEtaEllEtaEta} to estimate
\begin{align*}
 \|d_{\eta}\mathbf{E}(\eta)\dot{\eta} \|_{H^{s-\f32-\delta}} &\leq \|d_{\eta}R_E(\eta)\dot{\eta} \|_{H^{s-\f32-\delta}} + \|T_{\ell}\dot{\eta}\|_{H^{s-\f32-\delta}} + \|T_{d_{\eta}\ell\dot{\eta}}\eta\|_{H^{s-\f32-\delta}} \\
 &\lesssim_{\|\eta\|_{H^s}} \|\eta\|_{H^{s+\f52}} \| \dot{\eta}\|_{H^s} + \| \eta\|_{H^s}\|\dot{\eta}\|_{H^{s+\f52-\delta}}.
\end{align*}
This leads to the estimate \eqref{ElasticDiff}.
\end{proof}

\subsection{Contraction estimate for one-phase Muskat problem}
Suppose that $\eta_1$ and $\eta_2$ are two solutions of one-phase Muskat problem \eqref{OneMuskat} in $Z^s(T)$ with the condition \eqref{BottomH}, then we show that the $Z^s(T)$ norm of $\eta_1 - \eta_2$ can be controlled by the initial condition of $\eta_1 - \eta_2$ in $H^s$.
\begin{proposition}\label{t:propConOne}
Let $s>2$.
Suppose that $\eta_1$ and $\eta_2$ are two solutions of one-phase Muskat problem \eqref{OneMuskat} in $Z^s$ with the condition \eqref{BottomH}, then
\begin{equation} \label{ContractOneMuskat}
 \|\eta_1 - \eta_2\|_{Z^s(T)} \lesssim_{\|(\eta_1, \eta_2)\|_{Z^s(T) \times Z^s(T)}} \|(\eta_1 - \eta_2)(0,\cdot) \|_{H^s}.
\end{equation}
The implicit constant in the inequality depends only on $\big(s,h, \frac{\sigma}{\mu^-},  \frac{\rho^- g}{ \mu^-}\big)$.
\end{proposition}

\begin{proof}
Writing $\delta\eta = \eta_1 - \eta_2$ for the difference of two solutions to the one-phase Muskat problem.
From the equation \eqref{OneMuskat}, one can write
\begin{equation} \label{DiffMuskatOne}
    \p_t \delta\eta = -\f{\sigma}{\mu^-}G^-(\eta_1)(\mathbf{E}(\eta_1) - \mathbf{E}(\eta_2))- \mathcal{R}_0,
\end{equation}
where the remainder term $\mathcal{R}_0$ is given by
\begin{equation*}
 \mathcal{R}_0 : = \f{\rho^- g}{\mu^-}G^-(\eta_1) \delta\eta + [G^-(\eta_1)-G^-(\eta_2)]\cdot \left( \f{\sigma}{\mu^-}\mathbf{E}(\eta_2) + \f{\rho^- g}{\mu^-}\eta_2\right).
\end{equation*}
Using the estimate \eqref{GfEtaf},  
\begin{equation*}
 \|G^-(\eta_1) \delta\eta \|_{H^{s-\f52}} \lesssim_{\|\eta_1\|_{H^s}} \|\delta \eta\|_{H^{s-\f32}}.
\end{equation*}
According to the difference estimate for the Dirichlet-Neumann operator \eqref{GetaDiff},  we estimate
\begin{align*}
& \left\| [G^-(\eta_1)-G^-(\eta_2)]\cdot \left( \f{\sigma}{\mu^-}\mathbf{E}(\eta_2) + \f{\rho^- g}{\mu^-}\eta_2\right) \delta\eta  \right\|_{H^{s-\f52}} \\
 &\lesssim_{\|(\eta_1, \eta_2)\|_{H^s \times H^s}} \left(  \f{\sigma}{\mu^-} +  \f{\rho^- g}{\mu^-}\right) \|\delta \eta\|_{H^s} \|\eta_2\|_{H^{s+\f52}}
\end{align*}
Next, we rewrite the term $G^-(\eta_1)(\mathbf{E}(\eta_1) - \mathbf{E}(\eta_2))$.
Using the paralinearization of the Dirichlet-Neumann operator \eqref{GfDRf},
\begin{equation*}
 G^-(\eta_1)(\mathbf{E}(\eta_1) - \mathbf{E}(\eta_2)) = |D|(\mathbf{E}(\eta_1) - \mathbf{E}(\eta_2)) + \mathcal{R}_1, \quad \|\mathcal{R}_2\|_{H^{s-\f52}}\lesssim_{\|\eta_1\|_{H^s}} \|\mathbf{E}(\eta_1) - \mathbf{E}(\eta_2) \|_{H^{s-\f32 - \delta}}.
\end{equation*}
For the difference estimate of two elastic terms, we recall the estimate in \eqref{ElasticDiff}, 
\begin{equation*}
\|\mathbf{E}(\eta_1) - \mathbf{E}(\eta_2) \|_{H^{s-\f32 - \delta}} \lesssim_{\|(\eta_1, \eta_2)\|_{H^s \times H^s}} \|\delta\eta\|_{H^{s+\f52 -\delta}} +\Big(\|\eta_1\|_{H^{s+\f52-\delta}}+\|\eta_2\|_{H^{s+\f52-\delta}}\Big)\|\delta\eta\|_{H^{s}}.
\end{equation*}
Hence, we get the estimate for $\mathcal{R}_1$:
\begin{equation*}
\|\mathcal{R}_1\|_{H^{s-\f52}}\lesssim_{\|(\eta_1, \eta_2)\|_{H^s \times H^s}} \|\delta\eta\|_{H^{s+\f52 -\delta}}.
\end{equation*}
We then write 
\begin{equation*}
 |D|(\mathbf{E}(\eta_1) - \mathbf{E}(\eta_2)) = |D|T_{\ell_1}\delta\eta + |D|T_{\ell_1 - \ell_2}\eta_2 + |D|(R_E(\eta_1)-R_E(\eta_2)).
\end{equation*}
By the result in Lemma \ref{t:REEst}, 
\begin{equation*}
 \||D|(R_E(\eta_1)-R_E(\eta_2))\|_{H^{s-\f52}} \lesssim_{\|(\eta_1, \eta_2)\|_{H^s \times H^s}}  \Big(\|\eta_1\|_{H^{s+\f52}}+ \|\eta_2\|_{H^{s+\f52}} \Big) \|\delta \eta\|_{H^s}.
\end{equation*}
From \eqref{tdEtaEllEtaEta}, we get that
\begin{equation*}
 \| |D|T_{\ell_1 - \ell_2}\eta_2\|_{H^{s-\f52}} \lesssim_{\|(\eta_1, \eta_2)\|_{H^s \times H^s}}  \| \delta \eta\|_{H^s}\|\eta_2\|_{H^{s+\f52}}.
\end{equation*}
Collecting the above estimates, the equation \eqref{DiffMuskatOne} can be written as
\begin{equation} \label{ParaDiffMuskatOne}
\p_t\delta \eta=-\f{\sigma}{\mu^-}|D|T_{\ell_1}\delta\eta+\mathcal{R}_2,
\end{equation}
where the remainder term
\begin{equation*}
\|\mathcal{R}_2\|_{H^{s-\f52}}\lesssim_{\|(\eta_1, \eta_2)\|_{H^s \times H^s}} \|\delta\eta\|_{H^{s+\f52-\delta}} +\Big(\|\eta_1\|_{H^{s+\f52}}+\|\eta_2\|_{H^{s+\f52}}\Big)\| \delta \eta\|_{H^s}.
\end{equation*}
Note that the equation \eqref{ParaDiffMuskatOne} takes the similar form as the equation \eqref{ParaOneMuskat}.
Hence, doing the same analysis as in the proof of Proposition \ref{t:OnePhaseEst}, we get the energy inequality
\begin{align*}
\f12 \|\delta \eta \|^2_{H^s} &\leq - \f{1}{\mathcal{F}(\|(\eta_1, \eta_2)\|_{H^s\times H^s})} \|\delta \eta \|^2_{H^{s+\f52}} + \mathcal{F}(\|(\eta_1, \eta_2)\|_{H^s\times H^s}) \|\delta \eta\|_{H^{s+\f52}} \|\delta \eta\|_{H^{s+\f52-\delta}} \\
&+ \mathcal{F}(\|(\eta_1, \eta_2)\|_{H^s\times H^s})\|(\eta_1, \eta_2)\|_{H^{s+\f52}\times H^{s+\f52}}  \|\delta \eta\|_{H^s} \|\delta \eta\|_{H^{s+\f52}},
\end{align*}
for some function $\mathcal{F}: \R^+ \rightarrow \R^+$ depending only on $(h,s,\frac{\rho^- g}{\mu^-},\frac{\sigma}{\mu^-})$.
By interpolation and Young's inequality, we bound the last two terms in the inequality by 
\begin{align*}
&\mathcal{F}(\|(\eta_1, \eta_2)\|_{H^s\times H^s}) \|\delta \eta\|_{H^{s+\f52}} \|\delta \eta\|_{H^{s+\f52-\delta}} \\
\leq& \f{1}{4\mathcal{F}(\|(\eta_1, \eta_2)\|_{H^s\times H^s})} \|\delta \eta \|^2_{H^{s+\f52}} + \mathcal{F}_1(\|(\eta_1, \eta_2)\|_{H^s\times H^s}) \|\delta \eta\|^2_{H^s}, \\
& \mathcal{F}(\|(\eta_1, \eta_2)\|_{H^s\times H^s}) \Big(\|(\eta_1, \eta_2)\|_{H^{s+\f52}\times H^{s+\f52}}\Big)  \|\delta \eta\|_{H^s} \|\delta \eta\|_{H^{s+\f52}} \\
\leq& \f{1}{4\mathcal{F}(\|(\eta_1, \eta_2)\|_{H^s\times H^s})} \|\delta \eta \|^2_{H^{s+\f52}} + \mathcal{F}_1(\|(\eta_1, \eta_2)\|_{H^s\times H^s}) \|(\eta_1, \eta_2)\|_{H^{s+\f52}\times H^{s+\f52}}^2 \|\delta \eta\|^2_{H^s},
\end{align*}
for some functions $\mathcal{F}, \mathcal{F}_1$ and $\mathcal{F}_2$ depending only on $(h,s,\frac{\rho^- g}{\mu^-},\frac{\sigma}{\mu^-})$.
Hence, we get the inequality 
\begin{align*}
\f12 \|\delta \eta \|^2_{H^s} &\leq - \f{1}{\mathcal{F}(\|(\eta_1, \eta_2)\|_{H^s\times H^s})} \|\delta \eta \|^2_{H^{s+\f52}} + \mathcal{F}(\|(\eta_1, \eta_2)\|_{H^s\times H^s}) \|(\eta_1, \eta_2)\|_{H^{s+\f52}\times H^{s+\f52}}^2  \|\delta \eta\|_{H^s}.
\end{align*}
Applying the Gronwall's inequality and using the definition of the $Z^s(T)$ norm, we obtain the desired estimate \eqref{ContractOneMuskat}.
\end{proof}

\subsection{Contraction estimate for two-phase Muskat problem}
Suppose that $\eta_1$ and $\eta_2$ are two solutions of two-phase Muskat problem \eqref{TwoMuskatOne}-\eqref{TwoMuskatTwo} in $Z^s(T)$ with the condition \eqref{BottomH}, then similar to the one-phase case, we show that the $Z^s(T)$ norm of $\eta_1 - \eta_2$ can be controlled by the initial condition of $\eta_1 - \eta_2$ in $H^s$.

For $j = 1,2$, consider $f_j^\pm$ that solve the system 
\begin{equation*} 
\begin{cases}
f^-_j - f^+_j =  \sigma \mathbf{E}(\eta_j) + g(\rho^- - \rho^+)\eta_j, \\
\frac{1}{\mu^+}G^+(\eta_j)f^+_j = \frac{1}{\mu^-}G^-_j(\eta_j)f^-.
\end{cases}
\end{equation*}
We write $\delta f^\pm = f^\pm_1 - f^\pm_2$, and $\delta \eta = \eta_1 - \eta_2$ for the difference of these functions.
We first prove the following contraction estimate for $f^\pm$.

\begin{lemma}
Let $s\geq 2$, and $\delta \in (0, \f12]$.
For each $r\in [\f12, s-\f32]$,
\begin{equation} \label{DeltaFpmEst}
\begin{aligned}
 \|\delta f^\pm \|_{\tilde{H}^r_\pm} &\lesssim_{\|(\eta_1, \eta_2)\|_{H^s\times H^s}} \sigma \|\delta \eta \|_{H^{r+4}}+ (\rho^- -\rho^+)g \|\delta \eta\|_{H^r}\\
 +& (\sigma+(\rho^--\rho^+)g)\| \delta \eta\|_{H^s}\|(\eta_1, \eta_2)\|_{H^{s+\f52}\times H^{s+\f52}}.
\end{aligned}    
\end{equation}
The implicit constant in the inequality depends only on $\big(s,h, r,\mu^\pm\big)$.
\end{lemma}

\begin{proof}
We write
\begin{align*}
&\delta f^- - \delta f^+ = \delta k = \sigma (\mathbf{E}(\eta_1)-\mathbf{E}(\eta_2)) + g(\rho^- - \rho^+)\delta \eta, \\
&\f{1}{\mu^-}G^-(\eta_1)\delta f^- - \f{1}{\mu^+}G^+(\eta_1)\delta f^+ = \f{1}{\mu^+}[G^+(\eta_1)-G^+(\eta_2)]f_2^+ - \f{1}{\mu^-}[G^-(\eta_1)-G^-(\eta_2)]f_2^-.
\end{align*}
Applying $G^+(\eta_1)$ to the first equation, and using the fact that $G^\pm(\eta_1)\delta f^\pm = \mp |D|\delta f^\pm + R^\pm(\eta_1)\delta f^\pm$, we get
\begin{align*}
&|D|\delta f^- = \f{\mu^-}{\mu^+ + \mu^-}(\sigma|D|[\mathbf{E}(\eta_1)-\mathbf{E}(\eta_2)] + g(\rho^- - \rho^+)|D|\delta \eta) + \f{\mu^+ \mu^-}{\mu^++\mu^-}F, \\
F =& \f{1}{\mu^+}R^+(\eta_1)\delta f^+ -  \f{1}{\mu^-}R^-(\eta_1)\delta f^- + \f{1}{\mu^+}[G^+(\eta_1)-G^+(\eta_2)]f^+_2  - \f{1}{\mu^-}[G^-(\eta_1)-G^-(\eta_2)]f_2^-.
\end{align*}
Using the contraction estimate \eqref{ElasticDiff}, 
\begin{equation*}
 \|\delta f^-\|_{H^r}  \lesssim_{\|(\eta_1, \eta_2)\|_{H^{s}\times H^{s}}} \sigma \|\delta \eta \|_{H^{r+4}}+ (\rho^- -\rho^+)g \|\delta \eta\|_{H^r} + \|F \|_{H^{r-1}} + \sigma \| (\eta_1,\eta_2) \|_{H^{r+4}\times H^{r+4}}\| \delta \eta\|_{H^{s}} .
\end{equation*}
Using estimates \eqref{GfDRf}, \eqref{GetaDiff} and \eqref{fpmEstR},
\begin{align*} 
&\| R^\pm(\eta_1)\delta f^\pm\|_{H^{r-1}} \lesssim_{\|(\eta_1,\eta_2)\|_{H^s\times H^s}}\|\delta f^{\pm}\|_{\tilde{H}^{r-\delta}_{\pm}}, \\
 &\|[G^{\pm}(\eta_1) - G^{\pm}(\eta_2)]f_2^{\pm}\|_{H^{s-\f52}} \lesssim_{\|(\eta_1, \eta_2)\|_{H^s \times H^s}} (\sigma + (\rho^- -\rho^+)g)\|\delta \eta\|_{H^s} \|(\eta_1, \eta_2)\|_{H^{s+\f52}\times H^{s+\f52}}.
\end{align*}
For $r\in[\f12,s-\f32]$, we then have
\begin{align*}
\|F\|_{H^{r-1}}
\lesssim_{\|(\eta_1,\eta_2)\|_{H^s\times H^s}}\|\delta f^{\pm}\|_{\tilde{H}^{r-\delta}_{\pm}}
+(\sigma + (\rho^- -\rho^+)g)\|\delta \eta\|_{H^s} \|(\eta_1, \eta_2)\|_{H^{s+\f52}\times H^{s+\f52}}.
\end{align*}
Note that $\delta f^+=\delta f^--\delta k$, we then obtain that
\begin{align*}
\|\delta f^{\pm}\|_{H^r}  \lesssim_{\|(\eta_1, \eta_2)\|_{H^{s}\times H^{s}}}& \|\delta f^{\pm}\|_{\tilde{H}^{r-\delta}_{\pm}}+\sigma \|\delta \eta \|_{H^{r+4}}+ (\rho^- -\rho^+)g \|\delta \eta\|_{H^r}\\
&+ (\sigma + (\rho^- -\rho^+)g)\| (\eta_1,\eta_2) \|_{H^{s+\f52}\times H^{s+\f52}}\| \delta \eta\|_{H^s} .
\end{align*}
From the definition of $\tilde{H}^{r}$, following the proof of (4.22) in \cite{MR4131404}, we can obtain
\begin{align*}
&\|\delta f^{\pm}\|_{H^{1,r}}\lesssim_{\|(\eta_1,\eta_2)\|_{H^s\times H^s}}\|\delta f^{\pm}\|_{\tilde{H}^{\f12}_{\pm}}+\|\delta f^{\pm}\|_{H^{1,r-\delta}}\\
&+\sigma \|\delta \eta \|_{H^{r+4}}+ (\rho^- -\rho^+)g \|\delta \eta\|_{H^r}+ (\sigma + (\rho^- -\rho^+)g)\| (\eta_1,\eta_2) \|_{H^{s+\f52}\times H^{s+\f52}}\| \delta \eta\|_{H^s}.
\end{align*}

By using the variational form the definition for two-phase Muskat problem in \cite{MR4090462} and the definition of $\tilde{H}^{\f12}_{\pm}$, we have
\begin{align*}
&\|\delta f^{\pm}\|_{\tilde{H}^{\f12}_{\pm}}  \lesssim_{\|(\eta_1, \eta_2)\|_{H^{s}\times H^{s}}} \|\mathbf{E}(\eta_1)-\mathbf{E}(\eta_2)\|_{H^{\f12}}+ \|\delta\eta\|_{H^s}(\|\mathbf{E}(\eta_1)\|_{H^{\f12}}+\|\mathbf{E}(\eta_2)\|_{H^{\f12}})\\
\lesssim&_{\|(\eta_1, \eta_2)\|_{H^{s}\times H^{s}}}\sigma\|\delta\eta\|_{\f92}+(\rho^--\rho^+)g\|\delta\eta\|_{H^{\f12}}+(\sigma+(\rho^--\rho^+)g)\|\delta\eta\|_{H^s}\left(\|\eta_1\|_{H^{\f92}}+\|\eta_2\|_{H^{\f92}}\right).
\end{align*}
From the above estimates, by using an induction argument, we  obtain \eqref{DeltaFpmEst}.
\end{proof}

We then prove the contraction estimate for two solutions $\eta_1$ and $\eta_2$.
\begin{proposition} \label{t:propConTwo}
Let $s\geq 2$.
Suppose that $\eta_1$ and $\eta_2$ are two solutions of one-phase Muskat problem \eqref{TwoMuskatOne}-\eqref{TwoMuskatTwo} in $Z^s$ with the condition \eqref{BottomH}, then
\begin{equation*} %\label{ContractTwoMuskat}
 \|\eta_1 - \eta_2\|_{Z^s(T)} \lesssim_{\|(\eta_1, \eta_2)\|_{Z^s(T) \times Z^s(T)}} \|(\eta_1 - \eta_2)(0,\cdot) \|_{H^s}.
\end{equation*}
The implicit constant in the inequality depends only on $\big(s,h, \sigma, \mu^\pm, (\rho^- -\rho^+)g\big)$.
\end{proposition}

\begin{proof}
From the equation \eqref{TwoMuskatOne}, $\delta \eta$ solves the equation
\begin{equation*}
\p_t \delta\eta = -\f{1}{\mu^-}G^-(\eta_1)\delta f^- - \f{1}{\mu^-}[G^-(\eta_1) - G^-(\eta_2)]f_2^-.
\end{equation*}
Using estimates \eqref{GetaDiff} and \eqref{fpmEstR},
\begin{equation*}
 \|[G^-(\eta_1) - G^-(\eta_2)]f_2^- \|_{H^{s-\f52}} \lesssim_{\|(\eta_1, \eta_2)\|_{H^s \times H^s}} (\sigma + (\rho^- -\rho^+)g)\|\delta \eta\|_{H^s} \|(\eta_1, \eta_2)\|_{H^{s+\f52}\times H^{s+\f52}}.
\end{equation*}
Using the paralinearization of the Dirichlet-Neumann operator \eqref{GfDRf}, we write
\begin{equation*}
    G^-(\eta_1) \delta f^- = |D|\delta f^- + R^-(\eta_1)\delta f^-,
\end{equation*}
where the remainder term satisfies
\begin{align*}
 &\|R^-(\eta_1)\delta f^- \|_{H^{s-\f52}}\lesssim_{\|(\eta_1, \eta_2)\|_{H^s \times H^s}}\|\delta f^-\|_{\tilde{H}^{s-\f32-\delta}}\\
 &\lesssim_{\|(\eta_1, \eta_2)\|_{H^s\times H^s}} \sigma \|\delta \eta \|_{H^{s+\f52-\delta}}+ (\rho^- -\rho^+)g \|\delta \eta\|_{H^{s-\f32-\delta}}\\
 &+ (\sigma + (\rho^- -\rho^+)g)\|\delta \eta\|_{H^s} \|(\eta_1, \eta_2)\|_{H^{s+\f52}\times H^{s+\f52}},
\end{align*}
where we have used \eqref{DeltaFpmEst} with $r=s-\f32-\delta$.
We then get that
\begin{align*}
\p_t\delta\eta=-\f{1}{\mu^-}|D|\delta f^-+\mathcal{R}_1,
\end{align*}
where the remainder term satisfies
\begin{align*}
\|\mathcal{R}_1\|_{H^{s-\f52}}\lesssim_{\|(\eta_1, \eta_2)\|_{H^s\times H^s}}(\sigma+(\rho^--\rho^+)g) \left(\|\delta \eta \|_{H^{s+\f52-\delta}}
 + \| \delta \eta\|_{H^s}\|(\eta_1, \eta_2)\|_{H^{s+\f52}\times H^{s+\f52}}\right).
\end{align*}
From the computations of above lemma, 
\begin{align*}
&|D|\delta f^- = \f{\mu^-}{\mu^+ + \mu^-}(\sigma|D|[\mathbf{E}(\eta_1)-\mathbf{E}(\eta_2)] + g(\rho^- - \rho^+)|D|\delta \eta) + \f{\mu^+ \mu^-}{\mu^++\mu^-}F,
\end{align*}
where the remainder term satisfies
\begin{align*}
\|F\|_{H^{s-\f52}}\lesssim_{\|(\eta_1, \eta_2)\|_{H^s\times H^s}}(\sigma+(\rho^--\rho^+)g) \left(\|\delta \eta \|_{H^{s+\f52-\delta}}
 + \| \delta \eta\|_{H^s}\|(\eta_1, \eta_2)\|_{H^{s+\f52}\times H^{s+\f52}}\right).
\end{align*}
From the calculations of \eqref{ParaDiffMuskatOne}, we have
\begin{align*}
\||D|[\mathbf{E}(\eta_1)-\mathbf{E}(\eta_2)]-|D|T_{\ell_1}\delta\eta\|_{H^{s-\f52}}
\lesssim_{\|(\eta_1, \eta_2)\|_{H^s\times H^s}}\|\delta \eta \|_{H^{s+\f52-\delta}}
 + \| \delta \eta\|_{H^s}\|(\eta_1, \eta_2)\|_{H^{s+\f52}\times H^{s+\f52}}.
\end{align*}
Thus, we obtain that
\begin{equation} \label{ParaDiffMuskatTwo}
\p_t\delta \eta=-\f{\sigma}{\mu^-}|D|T_{\ell_1}\delta\eta+\mathcal{R}_2,
\end{equation}
where the remainder term $\mathcal{R}_2$ satisfies
\begin{align*}
\|\mathcal{R}_2\|_{H^{s-\f52}}\lesssim_{\|(\eta_1, \eta_2)\|_{H^s\times H^s}}(\sigma+(\rho^--\rho^+)g) \left(\|\delta \eta \|_{H^{s+\f52-\delta}}
 + \| \delta \eta\|_{H^s}\|(\eta_1, \eta_2)\|_{H^{s+\f52}\times H^{s+\f52}}\right).
\end{align*}
Since \eqref{ParaDiffMuskatTwo} takes the same form as \eqref{ParaDiffMuskatOne}, the result follows from an argument analogous to that in Proposition \ref{t:propConOne}.
\end{proof}

\subsection{Proof of local well-posedness}
We now finish the proof of local well-posedness by adapting the a priori estimates and contraction estimates we have shown.
The argument here is essentially the same as Section $3.4$ and $4.4$ in \cite{MR4131404}.
We assume that $\eta_0 \in H^s(\R)$ for $s\geq 2$ satisfying dist$(\eta_0, \Gamma^\pm)>2h$. 

We first begin with the one-phase Muskat problem.
Let $J_\varepsilon$ be the mollifier that selects the frequency portion not greater than $\varepsilon^{-1}$ of a function.
Then we consider the approximate solution $\eta_\varepsilon$ that solves the  following ODE in the Banach space 
\begin{equation*}
    \p_t \eta_\varepsilon = -\f{1}{\mu^-} J_\varepsilon [G^-(J_\varepsilon \eta_\varepsilon)(\sigma \mathbf{E}(J_\varepsilon\eta_\varepsilon)+ \rho^-gJ_\varepsilon\eta_\varepsilon)], \quad  \eta_\varepsilon|_{t = 0} = \eta(0, x).
\end{equation*}
The solution $\eta_\varepsilon$ exists on some maximal time interval $(0, T_\varepsilon]$. 
A priori estimates Proposition \ref{t:OnePhaseEst} and Lemma \ref{t:DistBound} also hold for $\eta_\varepsilon$.
Hence, using a continuity argument, there exists a positive time $T<T_\varepsilon$ for all $\varepsilon \in (0,1]$ such that
\begin{equation*}
  \| \eta_\varepsilon\|_{Z^s(T)} \lesssim \mathcal{F}\big(\|\eta(0,\cdot)\|_{H^s}\big), \quad  \inf_{t\in[0,T]} \text{dist}(\eta_\varepsilon(t), \Gamma^{-}) >h,
\end{equation*}
for some function $\mathcal{F}$ that depends on $\big(s,h, \frac{\sigma}{\mu^-},  \frac{\rho^- g}{ \mu^-}\big)$.
The contraction estimate Proposition \ref{t:propConOne} also holds for $\eta_\varepsilon$.
By choosing $\varepsilon \rightarrow 0$, we obtain that the limiting function $\eta$ is a solution of \eqref{OneMuskat} in $Z^s(T)$ with initial data $\eta_0$.
The uniqueness of the solution and Lipschitz dependence of initial data for the solution follows directly from Proposition \ref{t:propConOne}.

As for two-phase Muskat problem, we consider the ODE
\begin{equation*}
    \p_t \eta_\varepsilon = -\f{1}{\mu^-} J_\varepsilon [G^-(J_\varepsilon \eta_\varepsilon)(J_\varepsilon f_\varepsilon^-)], \quad  \eta_\varepsilon|_{t = 0} = \eta(0, x),
\end{equation*}
where functions $f^\pm_\varepsilon$ solve
\begin{equation*}  
\begin{cases}
f^-_\varepsilon - f^+_\varepsilon = \sigma \mathbf{E}(\eta_\varepsilon) + g(\rho^- - \rho^+)\eta_\varepsilon, \\
\frac{1}{\mu^+}G^+(\eta_\varepsilon)f^+_\varepsilon = \frac{1}{\mu^-}G^-(\eta_\varepsilon)f^-_\varepsilon.
\end{cases}
\end{equation*}
The solvability and regularity of $f^\pm$ follow from Lemma \ref{t:fpmHHalfEst}.
A priori estimates Proposition \ref{t:TwoPhaseEst}, and the contraction estimate Proposition \ref{t:propConTwo} also hold for $\eta_\varepsilon$.
By passing to the limit $\varepsilon \rightarrow 0$, and doing the same as in the one-phase Muskat case, we obtain the local well-posedness of two-phase Muskat problem.


\section{Global well-posedness of Muskat problem with an elastic interface} \label{s:Global}
In this section, we prove Theorem \ref{t:global}, establishing the global well-posedness for the Muskat problem with an elastic interface under the assumption of small initial data.
This section transitions from the local energy-based arguments to an integral formulation, utilizing fixed-point lemmas to prove that solutions exist for all time $T>0$.

The overarching strategy in this section is to treat the nonlinear system as a perturbation of a fractional heat equation.
By rewriting the equations in integral form, we can apply space-time estimates of the fractional heat equation to show that the nonlinearities remain small if the initial data is sufficiently small.

\subsection{Global well-posedness for one-phase Muskat problem} \label{s:GlobalOneMuskat}
For one-phase Muskat problem, we rewrite the governing equation \eqref{OneMuskat} as:
\begin{equation*} 
\begin{aligned}
&\left(\p_t + \f{\sigma}{\mu^-}|D|^5 + \frac{\rho^- g}{\mu^-}|D| \right) \eta\\
=& -\f{\sigma}{\mu^-}R^-(\eta)\mathbf{E}(\eta) - \f{\sigma}{\mu^-} |D|(T_{\ell(\eta)}\eta+R_{E(\eta)}\eta - |D|^4\eta)  -\f{\rho^- g}{\mu^-} R^-(\eta) \eta.
\end{aligned}
\end{equation*}
Hence, $\eta$ is the solution to the integral equation
\begin{equation} \label{OneMuskatIntegral}
    \eta(t) = e^{- \f{\sigma}{\mu^-}t|D_x|^5 - \frac{\rho^- g}{\mu^-}t|D_x|} \eta_0 - \mathcal{B}(\eta, \eta),
\end{equation}
where the bilinear form is given by
\begin{equation*} %\label{OneMuskatBilinear}
\begin{aligned}
\mathcal{B}(\eta, f)(t,x) = \f{\sigma}{\mu^-}\int_0^t& e^{- \f{\sigma}{\mu^-}(t- \tau)|D_x|^5 - \frac{\rho^- g}{\mu^-}(t-\tau)|D_x|}\Big( R^-(\eta)(T_{\ell(\eta)} f + R_{E(\eta)} (f))\\
+&  |D|(T_{\ell(\eta)}f+R_{E(\eta)}f - |D|^4f)  +\f{\rho^- g}{\sigma} R^-(\eta) f \Big) d\tau.  
\end{aligned}
\end{equation*}

Let $X^s([0, T])$ denote the Banach space endowed with the equivalent norm
\begin{equation*}
 \|\eta\|_{X^s([0, T])}=\| \eta\|_{\wt L^\infty([0, T]; H^s)}+\f{\sigma}{\mu^-}\| \eta\|_{\wt L^1([0, T]; H^{s+5})}.
\end{equation*}
We prove the following result for the equation \eqref{OneMuskatIntegral}.
\begin{lemma} \label{t:OneMuskatFixed}
Let  $s > \f32$, there exists a small number $\delta>0$ such that if  $\| \eta_0\|_{H^s}<\delta$, then 
\begin{equation*}
    e^{- \f{\sigma}{\mu^-}t|D_x|^5 - \frac{\rho^- g}{\mu^-}t|D_x|} \eta_0 - \mathcal{B}(\eta, \eta),
\end{equation*} 
 has a unique fixed point $\eta$ in $X^s([0, T])$ for any $T>0$ with norm less than $C\| \eta_0\|_{H^s}$.
\end{lemma}
\begin{proof}
Let $T>0$, we apply the fixed point Lemma \ref{t:fixedpoint}. 
In the proof the constant $C$ may change from line to line.
From \eqref{SpaceTimeOne}, we have
\begin{align*}
\| e^{- \f{\sigma}{\mu^-}t|D_x|^5 - \frac{\rho^- g}{\mu^-}t|D_x|} \eta_0\|_{X^s([0,T])}
\leq C\|\eta_0\|_{H^s}.
\end{align*}
To match the paralinearization of the elastic term in \eqref{ElasticPara}, we define 
\begin{equation*}
 \mathbf{E}(\eta, f) : = T_{\ell(\eta)} f + R_{E(\eta)} (f).
\end{equation*}
Using \eqref{SpaceTimeOne} with $(q_1,q_2)=(\infty,1)$ and $(q_1,q_2)=(1,1)$, we have
\begin{align*}
\|\mathcal{B}(\eta,f)\|_{X^s([0,T])}
\leq& C\f{\sigma}{\mu^-}\|R^-(\eta)\mathbf{E}(\eta,f) \|_{\tilde{L}^1([0,T];H^s)}+C\f{\rho^- g}{\mu^-} \|R^-(\eta) f \|_{\tilde{L}^1([0,T];H^s)}\\
&+C\f{\sigma}{\mu^-}\||D|(T_{\ell(\eta)}f+R_{E(\eta)}f - |D|^4f)  \|_{\tilde{L}^1([0,T];H^s)}.
\end{align*}
We now estimate each term on the right-hand side of the inequality.
It follows from Lemma \ref{t:elliptic} that
\begin{align*}
\|R^-(\eta) f \|_{H^s}\leq C\|\eta\|_{W^{1+\e,\infty}}\|f\|_{H^{s+1}}+C\|\eta\|_{H^{s+1}}\|f\|_{W^{1,\infty}}.
\end{align*}
We also get that
\begin{align*}
&\|R^-(\eta)\mathbf{E}(\eta,f) \|_{H^s}
\leq C\|\eta\|_{W^{1+\e,\infty}}\|\mathbf{E}(\eta,f)\|_{H^{s+1}}+C\|\eta\|_{H^{s+1}}\|\mathbf{E}(\eta,f)\|_{H^s}\\
&\leq C\|\eta\|_{H^s}\|f\|_{H^{s+5}}+C\|f\|_{H^s}\|\eta\|_{H^{s+5}}
+C\|\eta\|_{H^{s+1}}\|f\|_{H^{s+4}}\\
&\leq C\|\eta\|_{H^s}\|f\|_{H^{s+5}}+C\|f\|_{H^s}\|\eta\|_{H^{s+5}}.
\end{align*}
Due to \eqref{ellEst}, we have 
\begin{align*}
\||D|(T_{\ell(\eta)}f+R_{E(\eta)}f - |D|^4f)\|_{H^s}\leq C\|\eta\|_{H^s}\|f\|_{H^{s+5}}+C\|f\|_{H^s}\|\eta\|_{H^{s+5}}.
\end{align*}
We then obtain that 
\begin{equation} \label{check:cB1}
\begin{aligned}
\|\mathcal{B}(\eta,f)\|_{X^s([0,T])}\leq&C\|\eta\|_{L^{\infty}([0,T];H^s)}\|f\|_{\tilde{L}^{1}([0,T];H^{s+5})}+C\|f\|_{L^{\infty}([0,T];H^s)}\|\eta\|_{\tilde{L}^{1}([0,T];H^{s+5})}\\
\leq&C\|\eta\|_{X^s([0,T])}\|f\|_{X^s([0,T])}.
\end{aligned}  
\end{equation}
From the contraction estimate \eqref{GetaDiff},
we can obtain
\begin{align*}
&\|[R^-(\eta_1)-R^-(\eta_2)]\mathbf{E}(\eta_2,f) \|_{H^s}\\
\leq& C\|\delta\eta\|_{H^s}[(\|\eta_1\|_{H^s}+\|\eta_2\|_{H^s})\|f\|_{H^{s+5}}+(\|\eta_1\|_{H^{s+5}}+\|\eta_2\|_{H^{s+5}})\|f\|_{H^{s}}],
\end{align*}
and by Lemma \ref{t:elliptic},
\begin{align*}
&\|R^-(\eta_1)(\mathbf{E}(\eta_1,f) -\mathbf{E}(\eta_2,f))\|_{H^s}\\
\leq& C(\|\eta_1\|_{H^s}+\|\eta_2\|_{H^s})\|\delta \eta\|_{H^{s+5}}\|f\|_{H^s}+C\|\delta\eta\|_{H^{s+5}}\|f\|_{H^{s}}.
\end{align*}
Similarly, we have
\begin{align*}
&\||D|(T_{\ell(\eta_1)}f+R_{E(\eta_1)}f - |D|^4f)-|D|(T_{\ell(\eta_2)}f+R_{E(\eta_2)}f - |D|^4f)\|_{H^s} \\
\leq& C\|\delta\eta\|_{H^s}[(\|\eta_1\|_{H^s}+\|\eta_2\|_{H^s})\|f\|_{H^{s+5}}+(\|\eta_1\|_{H^{s+5}}+\|\eta_2\|_{H^{s+5}})\|f\|_{H^{s}}].
\end{align*}
Thus, we get the difference estimate
\begin{equation} \label{check:cB2}
\begin{aligned}
\|\mathcal{B}(\eta_1,f)-\mathcal{B}(\eta_2,f)\|_{X^s([0,T])}\leq&(\|(\eta_1,\eta_2)\|_{X^s\times X^s([0,T])}+1)\|\delta \eta\|_{X^s([0,T])}\|f\|_{X^s([0,T])}.
\end{aligned}
\end{equation}
We then obtain this lemma by using the fixed-point Lemma \ref{t:fixedpoint}.
\end{proof}
Let $\eta_0^j\in H^s$, $j=1, 2$ be initial data with norm less than $\delta$ given in Lemma \ref{t:OneMuskatFixed}, and we denote $\eta^j$, $j = 1,2$ be the corresponding solutions with initial data $\eta_0^j$.
Using \eqref{OneMuskatIntegral}, \eqref{check:cB1} and \eqref{check:cB2}, we compute
\begin{align*}
 &\|\eta_1 - \eta_2\|_{X^s([0,T])} \leq \big\|e^{- \f{\sigma}{\mu^-}t|D_x|^5 - \frac{\rho^- g}{\mu^-}t|D_x|} (\eta_0^1 - \eta_0^2) \big\|_{X^s([0,T])} +  \|\mathcal{B}(\eta^1, \eta^1) - \mathcal{B}(\eta^2, \eta^2) \|_{X^s([0,T])} \\
 &\leq C\|\eta_0^1 - \eta_0^2 \|_{H^s} + \|\mathcal{B}(\eta^1, \eta^1- \eta^2) \|_{X^s([0,T])} +  \|\mathcal{B}(\eta^1, \eta^2) - \mathcal{B}(\eta^2, \eta^2) \|_{X^s([0,T])} \\
 &\leq C\|\eta_0^1 - \eta_0^2 \|_{H^s}   +C\|\eta^1\|_{X^s([0, T])}\| \eta^1-\eta^2\|_{X^s([0, T])}+C\| \eta^1-\eta^2\|_{X^s([0, T])}\| \eta^2\|_{X^s([0, T])}.
\end{align*}
Choosing $\delta$ small enough, we obtain that for any $T>0$,
\begin{equation*}
 \|\eta_1 - \eta_2\|_{X^s([0,T])} \lesssim \|\eta_0^1 - \eta_0^2 \|_{H^s}.
\end{equation*}
This shows the Lipschitz dependence on initial data of the solution map.
As a consequence, we obtain the global well-posedness for one-phase Muskat problem in $H^s$ in the sense of Hadamard.


\subsection{Global well-posedness for two-phase Muskat problem} \label{s:GlobalTwoMuskat}
For two-phase Muskat problem, we restrict ourself to the physically stable case where the denser fluid is on the bottom $\rho^+\leq \rho^-$.
One can write the system \eqref{TwoMuskatTwo} as 
\begin{equation*}
 \frac{1}{\mu^-}G^-(\eta)f^- = \frac{1}{\mu^+}G^+(\eta)f^- -\frac{g(\rho^- - \rho^+)}{\mu^+}G^+(\eta)\eta - \f{\sigma}{\mu^+}G^+(\eta) \mathbf{E}(\eta).
\end{equation*}
This can be further written as
\begin{equation*}
|D_x|f^- = \f{1}{\mu^+ + \mu^-}(\mu^- R^+(\eta)-\mu^+ R^-(\eta))f^-  -\f{g(\rho^- - \rho^+)\mu^-}{\mu^+ + \mu^-}G^+ (\eta) \eta -\f{\sigma \mu^-}{\mu^+ + \mu^-}G^+ (\eta) \mathbf{E}(\eta).
\end{equation*}
Hence, $f^-$ is the fixed-point of 
\begin{equation*} %\label{fminusFixed}
\begin{aligned}
\mathcal{K}(\eta)\varphi = &\f{1}{\mu^+ + \mu^-}(\mu^- |D_x|^{-1}R^+(\eta)-\mu^+ |D_x|^{-1}R^-(\eta))\varphi \\
-&\f{g(\rho^- - \rho^+)\mu^-}{\mu^+ + \mu^-}|D_x|^{-1}G^+ (\eta) \eta -\f{\sigma \mu^-}{\mu^+ + \mu^-}|D_x|^{-1}G^+ (\eta) \mathbf{E}(\eta). 
\end{aligned}
\end{equation*}

Using the equation \eqref{TwoMuskatOne}, we obtain that $\eta$ solves the nonlinear fractional heat equation.
\begin{equation*}  %\label{PDE:twophase}
\begin{aligned}
 &\p_t \eta + \frac{\sigma}{\mu^+ + \mu^-}|D_x|^5 \eta + \frac{g(\rho^- - \rho^+)}{\mu^+ + \mu^-}|D_x|\eta = -\frac{1}{\mu^+ + \mu^-}( R^+(\eta)+ R^-(\eta))f^- \\
 +& \frac{g(\rho^- - \rho^+)}{\mu^+ + \mu^-}R^+(\eta)\eta + \f{\sigma}{\mu^++\mu^-}R^+(\eta) \mathbf{E}(\eta) - \f{\sigma}{\mu^++\mu^-}|D_x|( \mathbf{E}(\eta) - \p_x^4\eta).
\end{aligned}
\end{equation*}

\begin{proposition}\label{Prop:f-}
Let $r>\f12$ and let $\eta\in H^s\cap H^{r+5}$, there exists a small number $\delta>0$ such that if  $\| \eta\|_{H^s}<\delta$, then the mapping $\mathcal{K}(\eta)$ has a unique fixed point $f^-$ in $H^s\cap H^{r+1}$ and
\begin{equation} \label{fMinusBound}
\|f^{-}\|_{H^{r+1}}\leq Cg(\rho^--\rho^+)\|\eta\|_{H^{r+1}}
+C\sigma\|\eta\|_{H^{r+5}}.
\end{equation}
\end{proposition}
\begin{proof}
We apply Lemma \ref{t:fp} with $E_1=H^s$, $E_2=H^{r+1}$ and 
\begin{align*}
u_0=-\f{g(\rho^- - \rho^+)\mu^-}{\mu^+ + \mu^-}|D_x|^{-1}G^+ (\eta) \eta -\f{\sigma \mu^-}{\mu^+ + \mu^-}|D_x|^{-1}G^+ (\eta) \mathbf{E}(\eta).
\end{align*}
From \eqref{GfDRf2}, we have
\begin{align*}
\|u_0\|_{H^{r+1}}\leq&g(\rho^--\rho^+)\||D_x|\eta\|_{H^{r}}
+\sigma\||D_x|\mathbf{E}(\eta)\|_{H^{r}}\\
\leq&Cg(\rho^--\rho^+)\|\eta\|_{H^{r+1}}
+C\sigma\|\eta\|_{H^{r+5}}.
\end{align*}
It remains to prove that the mapping
\begin{align*}
\mathcal{K}_1(\eta)\varphi := \f{1}{\mu^+ + \mu^-}(\mu^- |D_x|^{-1}R^+(\eta)-\mu^+ |D_x|^{-1}R^-(\eta))\varphi
\end{align*}
satisfies the condition in Lemma \ref{t:fp}. 
We apply the estimate \eqref{GfDRf2},
\begin{align*}
\|\mathcal{K}_1\varphi\|_{H^{s}}\leq C\|\varphi\|_{H^{s}},\quad
\|\mathcal{K}_1\varphi\|_{H^{r+1}}\leq C\|\eta\|_{H^s}\|\varphi\|_{H^{r+1}}+C\|\eta\|_{H^{r+1}}\|\varphi\|_{H^s}.
\end{align*}
Hence, there exists a unique fixed point $f^-\in H^s\cap H^{r+1}$ for $\mathcal{K}(\eta)$, and
\begin{align*}
\|f^-\|_{H^{r+1}}\leq Cg(\rho^--\rho^+)\|\eta\|_{H^{r+1}}
+C\sigma\|\eta\|_{H^{r+5}}+C\|\eta\|_{H^s}\|f^-\|_{H^{r+1}}+C\|\eta\|_{H^{r+1}}\|f^-\|_{H^s}.
\end{align*}
Taking $r=s-\f12$, then
\begin{align*}
\|f^-\|_{H^s}\leq Cg(\rho^--\rho^+)\|\eta\|_{H^s}
+C\sigma\|\eta\|_{H^{s+4}}.
\end{align*}
Thus, we obtain the estimate \eqref{fMinusBound}.
\end{proof}

Assume that $\|(\eta_1,\eta_2)\|_{H^s\times H^s}\leq\delta$, and $f^-_j$ is the fixed point of $\mathcal{K}(\eta_j)$. Then
\begin{align*}
\||D_x|\delta f^-\|_{H^{r}}&\lesssim
\|R^+(\eta_1)f^-_1-R^+(\eta_2)f^-_2\|_{H^r}+\|R^-(\eta_1)f^-_1-R^-(\eta_2)f^-_2\|_{H^r}\\
+&g(\rho^--\rho^+)\|G^+(\eta_1)\eta_1-G^+(\eta_2)\eta_2\|_{H^r}+\sigma\|G^+(\eta_1)\mathbf{E}(\eta_1)-G^+(\eta_2)\mathbf{E}(\eta_2)\|_{H^r}.
\end{align*}
Using \eqref{GetaDiff} and \eqref{GfDRf2}, we then obtain
\begin{equation}\label{Global-two-deltaf}
\begin{aligned}
\||D_x|\delta f^-\|_{H^{r}}\leq C(\sigma+g(\rho^--\rho^+))\left(\|(\eta_1, \eta_2)\|_{H^{r+5}\times H^{r+5}}\|\delta\eta\|_{H^s}+\|\delta\eta\|_{H^{r+5}}\right).
\end{aligned}   
\end{equation}

From Proposition \ref{Prop:f-} and \eqref{Global-two-deltaf}, performing the similar analysis as in the proof of global well-posedness for one-phase Muskat problem, we obtain the global well-posedness for two-phase Muskat problem.

\appendix
\section{Paradifferential and product estimates} \label{s:Paradifferential}
Here, we recall the definition of function spaces and some of the paradifferential estimates that are used in  previous sections.
Some of these results can be found in for instance \cite{MR2418072, MR2768550, MR3260858}.

\begin{definition}
We recall the  Littlewood-Paley frequency decomposition, $ I = \sum_{k\in \mathbb{N}} P_k$, where for each $k\geq 1$, $P_k$ are smooth symbols  localized at  frequency $2^k$,  and $P_0$ selects the low frequency components $|\xi|\leq 1$.
\begin{enumerate}
\item Let $s\in \mathbb{R}$, and $p,q \in [1, \infty]$.
The non-homogeneous Besov space $B^s_{p,q}(\mathbb{R})$ is defined as the space of all tempered distributions $u$ such that
\begin{equation*}
\| u\|_{B^s_{p,q}} : = \left\|(2^{ks}\|P_k u \|_{L^p})_{k=0}^\infty \right\|_{l^q} < +\infty.
\end{equation*}
\item When $p = q = \infty$, Besov space $B^s_{\infty, \infty}$ coincides with the \textit{Zygmund space} $C^s_{*}$.
When $p = q =2$, the Besov space $B^s_{2,2}$ becomes the \textit{Sobolev space} $H^s$.
\item One has the Sobolev embedding, 
\begin{equation}
 H^{s+\frac{1}{2}}(\mathbb{R}) \hookrightarrow C^s_{*}(\mathbb{R}) \quad \forall s, \label{HsCsEmbed}
\end{equation}
the Sobolev space $H^{s+\frac{1}{2}}(\mathbb{R})$ can be embedding into the Zygmund space $C^s_{*}(\mathbb{R})$.
\item Let $k\in \mathbb{N}$, we let $W^{k,\infty}(\mathbb{R})$ the space of all functions such that $\partial_x^j u \in L^\infty(\mathbb{R})$, $0\leq j \leq k$. 
For $\rho = k+ \gamma$ with $k\in \mathbb{N}$ and $\gamma \in (0,1)$, we denote $W^{\rho, \infty}(\mathbb{R})$  the
space of all function $u\in W^{k,\infty}(\mathbb{R})$ such that the $\partial_x^k u$ is $\gamma$- H\"{o}lder continuous on $\mathbb{R}$. 
\item The Zygmund space $C^s_{*}(\mathbb{R})$ is just the H\"{o}lder space $W^{s, \infty}(\mathbb{R})$ when $s\in (0,\infty)\backslash \mathbb{N}$.
One has the embedding properties
\begin{align*}
  &C_{*}^s(\mathbb{R}) \hookrightarrow L^\infty(\mathbb{R}), \quad s>0; \qquad L^\infty(\mathbb{R}) \hookrightarrow C_{*}^s, \quad s<0;\\
  &C_{*}^{s_1}(\mathbb{R})\hookrightarrow C_{*}^{s_2}(\mathbb{R}), \quad H^{s_1}(\mathbb{R})\hookrightarrow H^{s_2}(\mathbb{R}), \qquad s_1>s_2.
\end{align*}
\end{enumerate}
\end{definition}

\begin{definition}
\begin{enumerate}
\item Let $\rho\in [0,\infty)$, $m\in \mathbb{R}$. 
$\Gamma^m_\rho(\mathbb{R})$ denotes the space of locally bounded functions $a(x, \xi)$ on $\mathbb{R}\times (\mathbb{R}\backslash \{0\})$, which are $C^\infty$ with respect to $\xi$ for $\xi \neq 0$ and such that for all $k \in \mathbb{N}$ and $\xi \neq 0$, the function $x\mapsto \partial_\xi^k a(x,\xi)$ belongs to $W^{\rho,\infty}(\mathbb{R})$ and there exists a constant $C_k$ with
\begin{equation*}
\forall |\xi|\geq \frac{1}{2}, \quad \|\partial_\xi^k a(\cdot,\xi) \|_{W^{\rho,\infty}} \leq C_k (1+ |\xi|)^{m-k}.
\end{equation*}
Let $a\in \Gamma^m_\rho$,  we define the semi-norm
\begin{equation*}
M^m_{\rho}(a) = \sup_{k \leq \frac{3}{2}+\rho} \sup_{|\xi|\geq \frac{1}{2}} \|(1+ |\xi|)^{k-m}\partial_\xi^k a(\cdot,\xi)  \|_{W^{\rho,\infty}}.
\end{equation*}
\item Given $a\in \Gamma^m_\rho(\mathbb{R})$, let $C^\infty$ functions $\chi(\theta, \eta)$ and $\psi(\eta)$ be such that for some $0<\epsilon_1 < \epsilon_2<1$,
\begin{align*}
    &\chi(\theta, \eta) = 1,  \text{ if } |\theta| \leq \epsilon_1(1+ |\eta|), \qquad \chi(\theta, \eta) = 0,  \text{ if } |\theta| \geq \epsilon_2(1+ |\eta|),\\
    &\psi(\eta) = 0, \text{ if } |\eta|\leq \frac{1}{5}, \qquad \psi(\eta) =1, \text{ if } |\eta|\geq \frac{1}{4}.
\end{align*}
We define the paradifferential operator $T_a$ by
\begin{align*}
 \widehat{T_a u}(\xi) = \frac{1}{2\pi}\int \chi(\xi -\eta, \eta) \hat{a}(\xi-\eta, \eta)\psi(\eta)\hat{u}(\eta) d\eta,
\end{align*}
where $\hat{a}(\theta, \xi)$ is the Fourier transform of a with respect to the  variable $x$.

\item Let $m\in \mathbb{R}$, an operator  is said to be of order $m$ if, for all $s\in \mathbb{R}$, it is bounded from $H^s$ to $H^{s-m}$.
\end{enumerate}
\end{definition}

We recall the basic symbolic calculus for paradifferential operators in the following result.
\begin{lemma}[Symbolic calculus, \cite{MR2418072}]
Let $m\in \mathbb{R}$ and $\rho\in [0, +\infty)$.
\begin{enumerate}
\item If $a\in \Gamma^m_0$, then the paradifferential operator $T_a$ is of order m. Moreover, for all $s\in \mathbb{R}$, there exists a positive constant $K$ such that
\begin{equation}
\|T_a\|_{H^s\rightarrow H^{s-m}} \leq K M^m_0(a). \label{TABound}
\end{equation}
\item If $a\in \Gamma^m_\rho$, and $b\in \Gamma^{m^{'}}_\rho$ with $\rho>0$, then the operator $T_a T_b -T_{a\sharp b}$ is of order $m+m^{'}-\rho$, where the composition
\begin{equation*}
    a \sharp b := \sum_{\alpha<\rho} \frac{(-i)^\alpha}{\alpha !} \partial^\alpha_\xi a(x,\xi) \partial^\alpha_x b(x,\xi).
\end{equation*}
Moreover,  for all $s\in \mathbb{R}$, there exists a positive constant $K$ such that
\begin{equation}
  \|T_a T_b-T_{a\sharp b} \|_{H^s \rightarrow H^{s-m-m^{'}+\rho}} \leq K\left(M^m_\rho(a)M^{m^{'}}_0(b) + M^m_0(a)M^{m^{'}}_\rho(b)\right). \label{CompositionPara}
\end{equation}
\item Let $a\in \Gamma^m_\rho$ with $\rho > 0$. 
Denote by $(T_a)^{*}$ the adjoint operator of $T_a$ and by $\bar{a}$ the complex conjugate of $a$.
Then $(T_a)^{*} -T_{a^{*}}$ is of order $m - \rho$, where 
\begin{equation*}
    a^{*} =  \sum_{\alpha<\rho} \frac{1}{i^\alpha \alpha !} \partial^\alpha_\xi \partial^\alpha_x \bar{a}.
\end{equation*}
Moreover,  for all $s\in \mathbb{R}$, there exists a positive constant $K$ such that
\begin{equation}
 \|(T_a)^{*} -T_{a^{*}} \|_{H^s \rightarrow H^{s-m+\rho}} \leq KM^m_\rho(a). \label{AdjointBound}
\end{equation}
In particular, if $a$ is a function that is independent of $\xi$, then  $(T_a)^* = T_{\bar{a}}$.
\end{enumerate}
\end{lemma}

When $a$ is just a function, $T_a u$ becomes the low-high paraproduct.
Below, we record some estimates for products and paraproducts.

\begin{lemma}[\hspace{1sp}\cite{MR2768550}]
\begin{enumerate}
\item Let $s_0, s_1, s_2$ be such that $s_0 \leq s_2$ and $s_0< s_1 +s_2 -\frac{1}{2}$, then
\begin{equation}
    \|T_a u\|_{H^{s_0}(\mathbb{R})} \lesssim \| a\|_{H^{s_1}(\mathbb{R})} \|u \|_{H^{s_2}(\mathbb{R})}.
\end{equation}
If in addition to the conditions above, $s_1 +s_2 >0$, then
\begin{equation} \label{BalancedError}
\|au -T_a u \|_{H^{s_0}(\mathbb{R})} \lesssim \| u\|_{H^{s_1}(\mathbb{R})} \|a \|_{H^{s_2}(\mathbb{R})}.
\end{equation}
\item For $s>0$, then for $u, v\in H^s\cap L^\infty$, $uv\in H^s$, and 
\begin{equation} \label{HsLinftyEst}
 \|uv\|_{H^s} \lesssim \|u\|_{L^\infty} \|v\|_{H^s} + \|v\|_{L^\infty} \|u\|_{H^s}.
\end{equation}
\end{enumerate}
\end{lemma}

For nonlinear functions, we record below the Moser estimate, the difference estiamte and the paralinearization result.
\begin{lemma}[\hspace{1sp}\cite{MR2768550}]
\begin{enumerate}
\item (Moser) Let  a smooth function $F\in C^\infty(\mathbb{C}^N)$ satisfying $F(0) = 0$.
Then, there holds
\begin{equation} \label{MoserOne}
\|F(u) \|_{H^s} \lesssim_{\|u\|_{L^\infty}} \|u\|_{H^s},\quad s\geq 0.  
\end{equation}
\item  Let a smooth function $F\in C^\infty(\mathbb{C}^N)$ satisfying $\nabla F(0) = 0$.
For any $U,V \in H^s(\R^d)^N \cap L^\infty(\R^d)^N$, 
\begin{equation} \label{FUVDiffEst}
\begin{aligned}
& \|F(U)-F(V)\|_{H^s} \\
 \lesssim&_{\|U\|_{L^\infty}, \|V\|_{L^\infty}} \Big(  \|U-V\|_{H^s} + \|U-V\|_{L^\infty} \sup_{\tau\in [0,1]} \|V+\tau(U-V) \|_{H^s}\Big).
 \end{aligned}
\end{equation}
\item (Paralinearization) Let $s, \rho>0$, and $F(u)$ be a smooth function of $u$, then for any $u\in H^s(\mathbb{R}^d)\cap C^\rho_{*}(\mathbb{R}^d)$,
\begin{equation} \label{Paralinearization}
    \|F(u)-F(0)-T_{F^{'}(u)}u\|_{H^{s+\rho}(\mathbb{R}^d)} \leq C(\|u\|_{L^\infty(\mathbb{R}^d)})\|u\|_{C^\rho_{*}(\mathbb{R}^d)}\|u\|_{H^s(\mathbb{R}^d)}.
\end{equation}
\end{enumerate}
\end{lemma}

\section{Results on the Dirichlet-Neumann operator and parabolic estimates} \label{s:DNOperator}
In this section, we recall some results on the Dirichlet-Neumann operator.
We consider the following screened fractional Sobolev space defined in Leoni and Tice \cite{MR3989147}:
\begin{equation*}
 \widetilde{H}^\f{1}{2}_\Upsilon (\R) = \left\{ f\in \mathcal{S}'(\R) \cap L^2_{loc}(\R) : \int_\R \int_{B_{\R}(0, \Upsilon(x))} \dfrac{|f(x+y ) - f(x)|^2}{|y|^2} dydx < \infty  \right\} / \R, 
\end{equation*}
where $\Upsilon : \R \rightarrow (0, \infty]$ is a given lower semi-continuous function.
For the lower domain $\Omega^-$, we choose
\begin{equation*}
\Upsilon(x) = \begin{cases}
\infty, \quad \text{if }  \Gamma^- = \emptyset, \\
\mathfrak{d}_- (x): = \frac{\eta(x) - \underline{b}^-(x)}{2(\| \eta_x\|_{L^\infty} + \| \underline{b}^-_x\|_{L^\infty})}, \quad \text{if } \underline{b}^- \in \dot{W}^{1,\infty}(\R). 
\end{cases}
\end{equation*}
For the upper domain $\Omega^+$, $\underline{b}^-$ is replaced by $\underline{b}^+$. 
We define the space
\begin{equation*}
 \widetilde{H}^\f{1}{2}_\pm (\R) = \begin{cases}
\widetilde{H}^\f{1}{2}_\infty (\R), \quad \text{if }  \Gamma^- = \emptyset, \\
\widetilde{H}^\f{1}{2}_{\mathfrak{d}_\pm}(\R), \quad \text{if } \underline{b}^- \in \dot{W}^{1,\infty}(\R). 
\end{cases}
\end{equation*}

Let $\mathcal{S}'(\R)$ and and $\mathcal{P}(\R)$ denote the space of tempered distributions on $\R$ and the set of polynomials on $\R$.
We also define the slightly-homogeneous Sobolev spaces
\begin{equation*}
H^{1,\sigma}(\R) = \{f \in \mathcal{S}'(\R)\cap L^2_{loc}(\R): f_x \in H^{\sigma-1}(\R) \}/ \R .
\end{equation*}
For $s>\f12$, we set 
\begin{equation*}
   \widetilde{H}^s_\pm (\R) = \widetilde{H}^\f{1}{2}_\pm(\R)\cap H^{1,s}(\R). 
\end{equation*}

Using the Sobolev spaces defined above, we have the following result for the  Dirichlet-Neumann operator.

\begin{lemma} [\hspace{1sp}\cite{MR3260858, MR4090462}]
Let $s>\f32$, and $\f12 \leq \mathfrak{s} \leq s$.
Consider $f \in \widetilde{H}^\mathfrak{s}_- (\R)$ and $\eta \in H^s(\R)$ with dist$(\eta, \Gamma^-) \geq h>0$.
Then $G^-(\eta)f \in H^{\mathfrak{s}-1}(\R)$ and
\begin{equation} \label{GfEtaf}
 \|G^-(\eta)f \|_{H^{\mathfrak{s} -1}} \lesssim_{\|\eta\|_{H^s}}  \|f\|_{ \widetilde{H}^\mathfrak{s}_-},
\end{equation}
where the implicit constant in the inequality depends only on $s, \mathfrak{s}, h$ and $\| \underline{b}^-\|_{\dot{W}^{1,\infty}(\R)}$.
\end{lemma}

For the two-dimensional problem, due to the simple geometry, we have the following result on the paralinearization of the Dirichlet-Neumann operator.
\begin{lemma}[\hspace{1sp}\cite{MR3260858, MR4090462}]
Let $s>\f32$, $\delta \in (0, \f12]$, and $\mathfrak{s} \in [\f12, s-\delta]$.
If $f \in \widetilde{H}^\mathfrak{s}_- (\R)$ and $\eta \in H^s(\R)$ with dist$(\eta, \Gamma^-)>h>0$, then 
\begin{equation} \label{GfDRf}
G^-(\eta)f = |D|f + R^-(\eta)f, \quad \| R^-(\eta)f\|_{H^{\mathfrak{s} - 1+\delta}} \lesssim_{\|\eta\|_{H^s}} \|f\|_{\widetilde{H}^\mathfrak{s}_-},
\end{equation}
where the implicit constant in the inequality depend only on $(s, \mathfrak{s}, \delta, h)$. 
\end{lemma}

Finally, we record the contraction estimate for the Dirichlet-Neumann operator.
This will be needed when we prove the uniqueness and the continuous dependence of the solutions.
\begin{lemma}[\hspace{1sp}\cite{MR4090462}]
Let $s>\f32$.
If $f \in \widetilde{H}^{s-\f12}_- (\R)$ and $\eta_1, \eta_2 \in H^s(\R)$ with dist$(\eta_j, \Gamma^-)>h>0$ for $j = 1,2$, then for all $\mathfrak{s} \in [\f12, s]$
\begin{equation} \label{GetaDiff}
\| G^-(\eta_1)f - G^-(\eta_2)f\|_{H^{\mathfrak{s} - 1}} \lesssim_{\|(\eta_1, \eta_2)\|_{H^s \times H^s}}\|\eta_1 - \eta_2 \|_{H^s} \|f\|_{\widetilde{H}^{\mathfrak{s}}_-},
\end{equation}
where the implicit constant in the inequality depends only on $(s, \mathfrak{s}, h)$. 
\end{lemma}

Next, we obtain the space-time estimates for fractional heat equations. 
\begin{lemma}[\hspace{1sp}\cite{MR2768550}] \label{Heates1}
There exist positive constants $c$ and $C$, such that for all $p\in [1, \infty]$, $t>0$, $\alpha_1,\nu_1>0$, $\alpha_2,\nu_2\geq0$ and $k\in \mathbb{Z}$, we have
\begin{equation*}
 \|e^{-t(\nu_1|D_x|^{\alpha_1}+\nu_2|D_x|^{\al_2})}P_ju\|_{L^p(\R)} \leq C e^{-c (2^{\alpha_1j}+2^{\al_2j} )t} \|P_j u\|_{L^p(\R)}.
\end{equation*}
\end{lemma}
Note that our discussion is restricted to the case 
$p=2$, where the proof is quite simple.

By slightly modifying the proof of Proposition $2.9$ in \cite{MR4348695}, we obtain the following result for the space-time estimate for the fractional heat kernel.

\begin{proposition}
Let $s\in \R$, $\nu_1>0$, $\nu_2\geq0$, and $1\leq q_2 \leq q_1 \leq \infty$.
Let $\alpha_1>0$, $\alpha_2\geq0$, $I = [a,b]$, where $a\in \R \cup \{-\infty \}$ and $b\in \R$.
Then there exist positive constants $C_1$ and $C_2$, such that
\begin{align}
 &\|e^{-z(\nu_1  |D_x|^{\alpha_1}+\nu_2|D_x|^{\alpha_2})}u(x) \|_{\widetilde{L}^{q_1}_z(I; H^{s+\frac{\alpha_1}{q_1}})} \leq \frac{C_1}{\nu_1^{\frac{1}{q_1}}}\|u\|_{H^s},  \label{SpaceTimeOne}\\
 &\left\|\int_a^z e^{-(z-y)(\nu_1|D_x|^{\alpha_1}+\nu_2|D_x|^{\alpha_2})}f(x,y)dy \right\|_{\widetilde{L}^{q_1}_z(I; H^{s+\frac{\alpha_1}{q_1}})} \leq \frac{C_2}{\nu_1^{1+\frac{1}{q_1}-\frac{1}{q_2}}} \|f\|_{\widetilde{L}^{q_2}(I; H^{s-\alpha_1+\frac{\alpha_1}{q_2}})}. \label{SpaceTimeTwo}
\end{align}
\end{proposition}
\begin{proof}
By Lemma \ref{Heates1}, we have
\begin{align*}
\|P_je^{-z(\nu_1  |D_x|^{\alpha_1}+\nu_2|D_x|^{\alpha_2})}u\|_{L^2(\R)}\leq& C e^{-c (\nu_12^{\alpha_1 j}+\nu_22^{\al_2j} )z} \|P_j u\|_{L^2(\R)}\\
\leq& C e^{-c \nu_12^{\alpha_1 j}z} \|P_j u\|_{L^2(\R)}.
\end{align*}
We then  get \eqref{SpaceTimeOne} from
\begin{align*}
\left(\int_{a}^be^{-c \nu_12^{\alpha_1j}q_1z}dz\right)^{\f{1}{q_1}}
=(cq_1)^{-\f{1}{q_1}}\f{2^{-\f{\al_1}{q_1} j}}{\nu_1^{\f1{q_1}}}\left(\int_a^be^{-c \nu_12^{\alpha_1j}q_1z}d(c\nu_12^{\al_1}q_1z)\right)^{\f1{q_1}}\leq C_1\f{2^{-\f{\al_1}{q_1} j}}{\nu_1^{\f1{q_1}}}.
\end{align*}

By Lemma \ref{Heates1} and Young's inequality in $z$, we have
\begin{align*}
&\left\|P_j\int_a^z e^{-(z-y)(\nu_1|D_x|^{\alpha_1}+\nu_2|D_x|^{\alpha_2})}f(x,y)dy \right\|_{L^{q_1}_z(I; L^2(\R))}\\
\leq& C\left\|\int_a^b e^{-c (\nu_12^{\alpha_1 j}+\nu_22^{\al_2j} )(z-y)}\|P_jf(\cdot,y)\|_{L^2}dy \right\|_{L^{q_1}_z(I)}\\
\leq& C\left\|\int_a^b e^{-c\nu_12^{\alpha_1 j}(z-y)}\|P_jf(\cdot,y)\|_{L^2}dy \right\|_{L^{q_1}_z(I)} \leq \f{C_2}{\nu_1^{1+\f1{q_1}-\f{1}{q_2}}}2^{\al_1j\big(-1+\f1{q_2}-\f{1}{q_1}\big)}\|P_jf\|_{L^{q_2}_z(I;L^2)},
\end{align*}
which implies \eqref{SpaceTimeTwo} by summing $j$ from $0$ to $\infty$.
\end{proof}
These results are used with either $\alpha_1 = 5, \alpha_2 = 1$ or $\alpha_1 = 1, \alpha_2 = 0$.

To obtain solutions of equations as fixed points of the mappings, we record here two fixed-point results from \cite{MR4348695}.
\begin{lemma}[\hspace{1sp}\cite{MR4348695}] \label{t:fixedpoint}
Let $(E, \|\cdot\|)$ be a Banach space, and let $\nu >0$.
Denote by $B_\nu$ the closed ball of radius $\nu$ centered at $0$ in $E$. 
Assume that $\mathcal{B}:E\times E \rightarrow E$ and there exists $\mathcal{F}: \R^+ \rightarrow \R^+$ such that
\begin{equation*} %\label{NuFTwoNu}
\nu \mathcal{F}(2\nu) \leq \f12,
\end{equation*}
and the following two conditions hold:
\begin{itemize}
\item For all $x\in B_\nu$, $\mathcal{B}(x, \cdot)$ is linear and
\begin{equation*} %\label{BUniformBound}
 \|\mathcal{B}(x, y)\| \leq \mathcal{F}(\|x\|)\|x\| \|y\|, \quad \forall y\in B_\nu.
\end{equation*}
\item For all $x_1, x_2, y \in B_\nu$,
\begin{equation*} %\label{BDiffBound}
 \|\mathcal{B}(x_1, y) - \mathcal{B}(x_2, y) \| \leq \mathcal{F}(\|x_1\| + \|x_2\|)\|x_1 - x_2\| \|y\|.
\end{equation*}
Then there exists $\delta = \delta(\nu, \mathcal{F})>0$ small enough such for all $x_0 \in E$ with norm less than $\delta$, $x \mapsto x_0 + \mathcal{B}(x,x)$ has a unique fixed point $x_*$ in $B_\nu$ with $\|x_*\| \leq 2\|x_0\|$.
\end{itemize}
\end{lemma}

\begin{lemma}[\hspace{1sp}\cite{MR4348695}]\label{t:fp}
Let $E_1$ and $E_2$ be two norm spaces such that $E_1$ is complete. Assume that $E_2$ has the Fatou property: if $(u_n)$ is a bounded sequence in $E_2$ then there exist $u\in E_2$ and a subsequence $(u_{n_k})$ such that $u_{n_k}\to u$ in a sense weaker than norm convergence in $E_1$ and that 
\[
\| u\|_{E_2}\le C\liminf_{k\to \infty} \| u_{n_k}\|_{ E_2}.
\]
 Assume that $K:E_1\to E_1$ is a linear such that $K: E_1\cap E_2\to E_2$ and the following property holds.  There exist $(\alpha_1, \alpha_2)\in (0, 1)^2$ and $A>0$ such that for all $u\in E_1\cap E_2$, 
\begin{align*}
&\| K(u)\|_{E_1}\le \alpha_1 \| u\|_{E_1}, \\
&\|K(u)\|_{E_2}\le \alpha_2\| u\|_{E_2}+A \| u\|_{E_1}. 
\end{align*}
Then for any $u_0\in E_1\cap E_2$ there exists a unique fixed point $u_*\in E_1\cap E_2$ of the mapping $u\mapsto u_0+K(u)$.
\end{lemma}

Finally, following the idea in Proposition $3.4$ in \cite{MR4348695}, we prove a refined remainder estimate of the Dirichlet-Neumann operator at higher regularity.
\begin{lemma}\label{t:elliptic}
Let $\eta ,f\in  H^{1+r}(\R)$ for $r>\f12$.
Suppose the lower boundary $\Gamma^-$ is either empty or flat $(\underline{b}^{-}_x = 0)$.
There exists a constant $c_1<1$ such that if
$\|\eta\|_{W^{1+\epsilon,\infty}}<c_1$ for $\epsilon$ sufficiently small, then
\begin{equation} \label{GfDRf2}
G^-(\eta)f = |D|f + R^-(\eta)f, \quad \| R^-(\eta)f\|_{H^{r}} \lesssim \| \eta\|_{W^{1+\epsilon,\infty}} \|f\|_{H^{r+1}} + \| f\|_{W^{1,\infty}} \|\eta\|_{H^{r+1}}.
\end{equation}
\end{lemma}

\begin{proof}
 For simplicity, we assume that $\Gamma^- = \emptyset$, so that the domain $\Omega^-$ has infinite depth.
 The analysis for the flat lower boundary case is similar, and one just needs to replace $H^{1+r}$ by $\widetilde{H}^{1+r}_-$.
 We fix the time $t$, and consider the elliptic problem in $\Omega^-$:
 \begin{equation}\label{eq:elliptic}
\begin{cases}
\Delta_{x, y} \phi=0\quad\text{in}~\Omega^-,\\
\phi=f\quad\text{on}~\Sigma,\\
\nabla_{x, y}\phi\to 0\quad\text{as}~y\to -\infty.
\end{cases}
\end{equation}
Assume that $\phi$ is a smooth solution of \eqref{eq:elliptic}. For $J=(-\infty, 0)$, we straighten the free boundary  using the change of variables $\R\times J\ni (x, z)\mapsto (x, \varrho(x, z))\in \Omega^-$, where
\begin{equation*}%\label{def:rho}
\varrho(x, z)=z+H(x, z),\quad H(x, z)=e^{z|D_x|}\eta(x),\quad (x, z)\in \R\times J.
\end{equation*}
Clearly, $\varrho(x, 0)=\eta(x)$ and $\varrho(x,z)\to -\infty$ as $z\to -\infty$.
Using \eqref{SpaceTimeOne}, we get
\begin{equation} \label{HRhoBound}
 \|H \|_{\widetilde{L} ^1(J; H^{r+5})} \lesssim \| \eta\|_{ H^{r+4}}, \quad  \|H \|_{\widetilde{L} ^\infty(J; H^{r})} \lesssim \| \eta\|_{ H^{r}}.
\end{equation}

Since $\phi$ is harmonic in $\Omega^-$, a direct computation shows that $v(x, z)=\phi(x,\varrho(x, z))$ solves the elliptic equation
\begin{equation}\label{eq:Deltav}
(\p_z+|D_x|)(\p_z-|D_x|)v=\partial_zQ_a[v]+|D_x|Q_b[v]\quad\text{in} ~\R \times J,
\end{equation}
 where $Q_a[v]$ and $Q_b[v]$ are given by
\begin{equation*}%\label{Qab}
\begin{aligned}
&Q_a[v]=\partial_xH\cdot \partial_xv-\frac{|\partial_xH|^2-|D_x|H}{1+|D_x|H}\partial_zv,\\
&Q_b[v]=|D_x|^{-1}\partial_x\big(\partial_xH\partial_zv-|D_x|H\partial_xv\big).
\end{aligned}
\end{equation*}

From the equation \eqref{eq:Deltav}, we get that $v$ is a fixed point of the operator 
\begin{equation}\label{def:cT}
\begin{aligned}
\mathcal{T}[v](x, z)&=e^{z|D_x|}f(x)+\int_0^ze^{(z-z')|D_x|}Q_a[v](x, z')dz'\\
&+\int_0^ze^{(z-z')|D_x|}\int_{-\infty}^{z'}e^{-(z'-\tau)|D_x|}|D_x|\{Q_b[v](x, \tau)-Q_a[v](x, \tau)\}d\tau dz'.
\end{aligned}
\end{equation}
If $v$  is a fixed point of $\mathcal{T}$ \eqref{def:cT}, then 
 the Dirichlet-Neumann operator $G^-(\eta)f$ is defined by
\begin{equation}\label{def:R}
G^-(\eta)f=|D_x|f+R^-(\eta)f:=|D_x|f+\int_{-\infty}^0e^{\tau |D_x|}|D_x|\{Q_b[v](x, \tau)-Q_a[v](x, \tau)\}d\tau.
\end{equation}
Define the auxiliary function
\begin{equation*}
    w(x,y) : = \int_{-\infty}^ye^{-(y-\tau)|D_x|}|D_x|\{Q_b[v](x, \tau)-Q_a[v](x, \tau)\}d\tau,\quad y\le 0.
\end{equation*}
The estimate \eqref{SpaceTimeTwo} with $\nu_1 = \alpha_1 = 1$, and $\nu_2 = \alpha_2 = 0$ shows that
\begin{equation}\label{Xest:w10}
 \| w\|_{X^r(J)}\le C \| |D_x|(Q_b-Q_a)\|_{\widetilde{L}^1(J;  H^{r+4})}\lesssim \| (Q_a, Q_b)\|_{\widetilde{L} ^1(J; H^{r+5}\times H^{r+5})}.
\end{equation}
We can also write the operator $\mathcal{T}$ as 
\begin{equation*}%\label{def:opK}
\mathcal{T}[v](x, z)=e^{z|D_x|}f(x)+K[v](x, z):=e^{z|D_x|}f(x)+\int_0^ze^{(z-z')|D_x|}\{w(x, z')+Q_a(x, z')\}dz'.
\end{equation*}
We have the estimate
\begin{equation*}
 \||D_x| (K[v]) \|_{X^r(J)} \lesssim \| w\|_{\widetilde{L} ^1(J; H^{r+5})} + \|Q_a\|_{\widetilde{L} ^1(J; H^{r+5})} \lesssim \| (Q_a, Q_b)\|_{\widetilde{L} ^1(J; H^{r+5}\times H^{r+5})}.
\end{equation*}
We compute the $z-$partial derivative of $K[v]$:
\begin{equation*}
\label{form:dzK}
\p_z K[v](x, z)=\int_0^ze^{(z-z')|D_x|}|D_x|\{w(x, z')+Q_a(x, z')\}dz'+w(x, z)+Q_a(x, z),
\end{equation*}
so that by using estimates \eqref{SpaceTimeTwo} and \eqref{Xest:w10}, 
\begin{align*}
 \|\p_z K[v](x, z)\|_{\wt L^1(J; H^{r+5})}&\lesssim\| w\|_{\wt L^1(H^{r+5})}+\| Q_a\|_{\wt L^1(J; H^{r+5})}\lesssim \| (Q_a, Q_b)\|_{\wt L^1(J; H^{r+5}\times H^{r+5})}, \\
  \|\p_z K[v](x, z)\|_{\wt L^\infty(J; H^{r})}&\lesssim \| w\|_{\wt L^1(H^{r+1})}+\| Q_a\|_{\wt L^1(J; H^{r+1})} + \| w\|_{\wt L^\infty (J;H^{r})} + \|Q_a\|_{\wt L^\infty(J; H^{r})} \\
 &\lesssim \| (Q_a, Q_b)\|_{\wt L^1(J; H^{r+5}\times H^{r+5})} + \|Q_a\|_{ \wt L^\infty(J; H^{r})}.
\end{align*}
Hence, we get that
\begin{align*}
\|( |D_x| K[v],\p_z K[v])\|_{X^r(J) \times X^r(J)} &\lesssim \| (Q_a, Q_b)\|_{\wt L^1(J; H^{r+5}\times H^{r+5})} + \|Q_a\|_{ \wt L^\infty(J; H^{r})} \\
&\lesssim  \|Q_a\|_{X^r(J)} +  \| Q_b\|_{\wt L^1(J;  H^{r+5})}.
\end{align*}

To estimate $Q_a[v]$, we need to consider the term
\begin{align*}
    I =&  \frac{|\partial_xH|^2-|D_x|H}{1+|D_x|H}\partial_zv = I_1 + I_2+ I_3 \\
    =& |\partial_xH|^2 \partial_zv - |\partial_xH|^2F(|D_x|H) \partial_zv  - F(|D_x|H)\partial_zv, \quad F(x) = \frac{x}{1+x}.
\end{align*}
We estimate $I$ using \eqref{HsLinftyEst}, Moser estimate \eqref{MoserOne}, and the fact that $|F|\leq 1$
\begin{align*}
 \|&I_1\|_{X^r(J)} \lesssim \|\p_x H \|^2_{L^\infty(J; L^\infty)} \|\p_z v\|_{X^r(J)} + \| \p_x H \|_{X^r(J)} \| \p_x H \|_{L^\infty(J; L^\infty)}\|\p_z v\|_{L^\infty(J; L^\infty)} \\
 &\lesssim \|\eta_x\|^2_{L^\infty} \|\p_z v\|_{X^r(J)} + \|\eta_x\|_{L^\infty} \|\eta\|_{H^{r+4}} \|\p_z v\|_{L^\infty(J; L^\infty)}, \\
 \|&I_2\|_{X^r(J)} \lesssim  \|\p_x H \|^2_{L^\infty(J; L^\infty)} \|\p_z v\|_{X^r(J)} + \| \p_x H \|_{X^r(J)} \| \p_x H \|_{L^\infty(J; L^\infty)}\|\p_z v\|_{L^\infty(J; L^\infty)} \\
 &+ \|\p_x H \|^2_{L^\infty(J; L^\infty)} \|\p_z v\|_{L^\infty(J; L^\infty)}  \|F(|D_x|H) \|_{X^r(J)}\\
 &\lesssim \|\eta_x\|^2_{L^\infty} \|\p_z v\|_{X^r(J)} + \|\eta_x\|_{L^\infty} \|\eta\|_{H^{r+4}} \|\p_z v\|_{L^\infty(J; L^\infty)} + \|\eta_x\|^2_{L^\infty} \|\p_z v\|_{L^\infty(J; L^\infty)} \|\eta\|_{H^{r+5}}, \\
 \|&I_3\|_{X^r(J)} \lesssim \|F(|D_x|H) \|_{L^\infty(J; L^\infty)} \|\p_z v\|_{X^r(J)} + \|F(|D_x|H) \|_{X^r(J)} \|\p_z v\|_{L^\infty(J; L^\infty)} \\
 &\lesssim \|\eta\|_{W^{1+\epsilon,\infty}}  \|\p_z v\|_{X^r(J)} + \| \eta\|_{H^{r+5}}   \|\p_z v\|_{L^\infty(J; L^\infty)}. 
\end{align*}
Other terms in $Q_a$ and $Q_b$ are estimated similarly, so that since $\|\eta\|_{W^{1,\infty}}< 1$,  
\begin{align*}
 \|( |D_x| K[v],\p_z K[v])\|_{X^r(J) \times X^r(J)} \lesssim&\|Q_a\|_{X^r(J)} +  \| Q_b\|_{\wt L^1(J;  H^{r+5})} \\
 \lesssim& \|\eta\|_{W^{1+\epsilon,\infty}} \|\nabla_{x,z} v\|_{X^r(J)} + \| \eta\|_{H^{r+5}}   \|\nabla_{x,z}  v\|_{L^\infty(J; L^\infty)}.
\end{align*}
We define the space
\begin{align*}
 X_*(J)& = L^\infty(J; L^\infty) \cap \wt L^1(J;  W^{1,\infty}), \quad  \mathcal{V}_*=\{v\in \mathcal{D}'(\R\times J): (|D_x|v, \p_zv)\in X_*(J)\}/\R, \\
  \mathcal{V}^r&=\{v\in \mathcal{D}'(\R\times J): (|D_x|v, \p_zv)\in   X^r(J)\}/\R.
\end{align*}
Then one has the estimate 
\begin{equation*}
 \|K[v] \|_{ \mathcal{V}^r} \lesssim \|\eta\|_{W^{1+\epsilon,\infty}} \| v\|_{\mathcal{V}^r(J)} + \| \eta\|_{H^{r+5}}   \|  v\|_{\mathcal{V}_*(J)},
\end{equation*}
and similarly, one can get the estimate
\begin{equation*}
 \|K[v] \|_{ \mathcal{V}_*} \lesssim \|\eta\|_{W^{1+\epsilon,\infty}} \| v\|_{\mathcal{V}_*(J)}.
\end{equation*}

Hence, according to Lemma \ref{t:fp}, using $E_1 = \mathcal{V}_*$ and $E_2 = \mathcal{V}^r$, by choosing $c_1$ small enough, $\mathcal{T}$ has a unique fixed point $v \in \mathcal{V}_* \cap  \mathcal{V}^r$.
Since $v- K[v] = e^{z|D_x|}f$, we have
\begin{align*}
   &\|\nabla_{x,z} v\|_{\wt L^\infty (J; L^\infty)}\lesssim \|\nabla_{x,z}e^{z|D_x|}f\|_{_{\wt L^\infty (J; L^\infty)}} \lesssim \|f\|_{W^{1,\infty}}, \\
   &\|\nabla_{x,z} v\|_{\wt L^\infty (J; H^r)}\lesssim \|\nabla_{x,z}e^{z|D_x|}f\|_{_{\wt L^\infty (J; H^r)}} \lesssim \|f\|_{H^{r+1}}.
\end{align*}
Note that using the definition of $R^-(\eta) f$ in \eqref{def:R}, and the estimate \eqref{SpaceTimeTwo},
\begin{align*}
 &\|R^-(\eta)f\|_{H^r} \lesssim \|(Q_a[v], Q_b[v]) \|_{\wt L^\infty (J; H^{r})}\\
 \lesssim &\|\eta\|_{W^{1+\epsilon,\infty}} \|\nabla_{x,z} v\|_{\wt L^\infty (J; H^{r})} + \|\eta\|_{H^{r+1}}  \|\nabla_{x,z} v\|_{\wt L^\infty (J; L^\infty)}\\
 \lesssim &\|\eta\|_{W^{1+\epsilon,\infty}} \|f\|_{ H^{r+1}} + \|\eta\|_{H^{r+1}}  \|f\|_{W^{1,\infty}}.
\end{align*}
This gives the estimate for the remainder term of the Dirichlet-Neumann operator.
\end{proof}

\textbf{Acknowledgments.}
Jiaqi Yang is supported by National Natural Science Foundation of China under Grant: 12471225.

\bibliography{HWW}
\bibliographystyle{plain}
	
\end{document}
